% ============================================================================
%  第六章  进阶定制与扩展
% ============================================================================

\chapter{进阶定制与扩展}

前几章介绍了模板的"开箱即用"与内部原理。本章给出进阶定制的常用
手段：改颜色、换字体、调版式、改页眉页脚，以及加载第三方宏包
的注意事项。

\section{颜色定制}

\subsection{切换主题色}

主题色通过文档类选项 \texttt{color} 选择，五套预设的
\texttt{structurecolor} 如图~\figref{fig:colors} 所示：

\begin{figure}[htbp]
  \centering
  \begin{tikzpicture}[font=\small]
    \definecolor{swgreen}{RGB}{0,120,2}
    \definecolor{swcyan}{RGB}{31,186,190}
    \definecolor{swblue}{RGB}{60,113,183}
    \definecolor{swgray}{RGB}{150,150,150}
    \definecolor{swblack}{RGB}{0,0,0}
    \foreach \col/\name [count=\i] in {swgreen/绿 (green),
                                       swcyan/青 (cyan),
                                       swblue/蓝 (blue),
                                       swgray/灰 (gray),
                                       swblack/黑 (black)}{
      \draw[fill=\col, draw=black!40]
            ({0.15+(\i-1)*1.35},0) rectangle ({1.15+(\i-1)*1.35},1.2);
      \node[align=center, font=\scriptsize] at ({0.65+(\i-1)*1.35},-0.35)
            {\name};
    }
  \end{tikzpicture}
  \caption{五套预设主题色}\label{fig:colors}
\end{figure}

\subsection{自定义颜色}

若预设色不合意，可重新定义模板使用的颜色变量。文档类使用的核心
颜色有 \texttt{structurecolor}（结构色）、\texttt{main}、
\texttt{second}、\texttt{third}（三个模块色）。在导言区覆盖即可：

\begin{lstlisting}
\definecolor{structurecolor}{RGB}{90,60,180}   % 改为紫色系
\definecolor{main}{RGB}{80,0,120}
\definecolor{second}{RGB}{200,60,0}
\definecolor{third}{RGB}{0,120,160}
\end{lstlisting}

由于定理盒子、页眉、标题、目录全部由这些变量着色，只需改动几行
即可实现全局换肤。

\begin{remark}
也可以用 \texttt{xcolor} 的 \verb|\colorlet| 把主题色映射到
其他已定义颜色，例如 \verb|\colorlet{main}{structurecolor}|，
实现"主色跟随结构色"的联动。
\end{remark}

\section{字体定制}

\subsection{中文字体}

文档类提供 \texttt{chinesefont} 选项控制中文字体方案：
\texttt{ctexfont}（默认，跟随 ctex 的字体配置）、
\texttt{founder}（方正字体）、\texttt{nofont}（不设置，
留给用户自行配置）。在 \texttt{nofont} 模式下，可以用
\texttt{fontspec} 的 \verb|\setCJKmainfont| 等命令指定系统字体：

\begin{lstlisting}
\documentclass[lang=cn, chinesefont=nofont]{elegantbook}
\setCJKmainfont{Source Han Serif SC}   % 思源宋体
\setCJKsansfont{Source Han Sans SC}    % 思源黑体
\setCJKmonofont{FangSong}              % 等宽/楷体
\end{lstlisting}

\subsection{西文与数学字体}

西文字体由 \texttt{math} 选项间接影响：\texttt{cm}（默认，
Computer Modern 风格）、\texttt{newtx}（Times 风格，正文与数学
同步替换）、\texttt{mtpro2}（商业数学字体，需要已安装
mtpro2 宏包）。数学字体切换对公式密集的书籍影响显著，
建议统一后全局使用。

\section{版式定制}

\begin{description}
  \item[\texttt{mode=simple}] 简约模式：定理类环境退化为
        \texttt{amsthm} 经典样式（无彩色盒子），适合纯文字书籍；
  \item[\texttt{device=pad}] 平板模式：纸张改为 $7.5\times10$ 英寸、
        边距 $16\,\mathrm{mm}$，为触屏阅读优化；
  \item[\texttt{titlestyle=display}] 通栏式章节标题
        （默认 \texttt{hang} 为悬挂式）；
  \item[\texttt{toc=twocol}] 目录双栏排版（默认 \texttt{onecol}）；
  \item[\texttt{marginpar=margintrue}] 开启边注栏
        （左 5\,cm 右 7\,cm 的版心偏移）；
  \item[\texttt{result=noanswer}] 关闭习题答案环境
        （默认 \texttt{answer}，显示答案）。
\end{description}

\section{页眉页脚定制}

页眉页脚由 \texttt{fancyhdr} 实现。默认页眉在左/右页分别显示
书名与章节名。若需自定义，直接使用 fancyhdr 的命令：

\begin{lstlisting}
\fancyhf{}
\fancyhead[L]{\small\leftmark}      % 左页眉：当前章节
\fancyhead[R]{\small\rightmark}     % 右页眉：当前小节
\fancyfoot[C]{\small\thepage}       % 页脚居中页码
\end{lstlisting}

注意 \texttt{leftmark}/\texttt{rightmark} 是 LaTeX 自动维护的
"当前章/节标题"标记宏，与 \verb|\chapter| 的 mark 机制联动。

\section{加载第三方宏包}

ElegantBook 不限制用户加载其他宏包（本手册即额外加载了
\texttt{algorithm} 与 \texttt{algpseudocode}）。实践中的建议：

\begin{itemize}
  \item 先查阅本手册附录 A 确认选项覆盖范围，避免重复配置；
  \item 宏包与文档类有交互时（如字体、几何、标题类），
        注意加载顺序：用户的 \verb|\usepackage| 在文档类之后，
        通常能覆盖模板默认；
  \item 使用 \texttt{tcolorbox} 系宏包时，直接用
        \texttt{\textbackslash tcbset} 追加样式即可，不要重新
        \texttt{\textbackslash newtcolorbox} 定义同名环境。
\end{itemize}

\section{与 ctex 的关系（为什么不必显式加载）}

如第一章所述，本手册全程不显式使用 ctex。原因在于：
ElegantBook 在 \texttt{lang=cn} 时已在内部加载 ctex（并配置
\verb|scheme=plain| 等选项），自动处理中文标点、缩进与字体。
用户若再 \verb|\usepackage{ctex}|，既可能造成选项冲突，也容易
因字体配置覆盖而出现警告。需要深度控制中文排版时，
应优先使用文档类提供的 \texttt{chinesefont} 选项与
\texttt{fontspec} 命令，而不是重复加载 ctex 宏包。

\section{完整换肤示例}

把前面几节的知识串起来：定义一个"紫罗兰"主题并应用到全文。

\begin{lstlisting}
% 导言区（documentclass 之后）
\definecolor{structurecolor}{RGB}{90,60,180}   % 标题、页眉、目录
\definecolor{main}{RGB}{80,0,120}              % 定义、例题盒子
\definecolor{second}{RGB}{200,60,0}            % 定理、引理、证明
\definecolor{third}{RGB}{0,120,160}            % 命题、假设、结论

% 把页眉装饰线条也改为新结构色（可选）
\renewcommand{\headrule}{\color{structurecolor}
  \hrule width\headwidth height 0.6pt \vspace{1pt}}
\end{lstlisting}

由于模板内部所有着色都引用这四个颜色变量，上述几行即可让封面、
标题、目录、定理盒子、页眉页脚整体换肤。下面的定理盒子即用
默认蓝色主题渲染，作为对比参照：

\begin{theorem}[换肤效果]
颜色定制只需覆盖四个核心变量，无需改动任何环境定义。
\end{theorem}

\section{页眉页脚深度定制}

模板使用 \texttt{fancyhdr}，其核心概念是"样式 + 分区命令"：

\begin{lstlisting}
\pagestyle{fancy}                 % 启用 fancy 样式
\fancyhf{}                        % 清空默认页眉页脚
\fancyhead[L]{\small\leftmark}    % 左页眉：当前章名
\fancyhead[R]{\small\rightmark}   % 右页眉：当前节名
\fancyfoot[C]{\small\thepage}     % 页脚居中页码
\renewcommand{\headrulewidth}{0.6pt}   % 页眉横线粗细
\renewcommand{\footrulewidth}{0pt}     % 页脚横线
\end{lstlisting}

\verb|\leftmark| 与 \verb|\rightmark| 是 LaTeX 内置的
"标记机制"输出：\verb|\chapter| 在内部调用
\verb|\markboth{章名}{章名}|，\verb|\section| 调用
\verb|\markright{节名}|，页眉据此显示"当前正在阅读哪一节"。
标记随内容自动更新，这是页眉与手写固定文字的
本质区别。

\section{章节标题样式定制}

标题样式由 \texttt{titlesec} 控制，可用
\verb|\titleformat| 修改：

\begin{lstlisting}
\titleformat{\section}
  {\Large\bfseries\color{structurecolor}}  % 格式
  {\thesection}                             % 编号
  {0.5em}                                   % 编号与标题间距
  {}                                        % 标题前内容
\titleformat{\chapter}[display]
  {\Huge\bfseries\color{structurecolor}}
  {\chaptername\ \thechapter}{20pt}{}
\end{lstlisting}

\verb|\titleformat| 的五个参数依次为：命令名、标题字体格式、
编号格式、编号与标题间的水平间距、标题文本前插入的内容；
可选参数 \texttt{[display]} 表示章标题独占多行（通栏式）。

\section{目录定制}

目录由 \texttt{tocloft} 控制，常用定制包括：

\begin{lstlisting}
% 章标题加粗着色
\renewcommand{\cftchapfont}{\bfseries\color{structurecolor}}
% 页码框宽度（页码右对齐）
\setlength{\cftchapnumwidth}{2em}
% 去除目录条目的点线（默认保留）
\renewcommand{\cftdot}{.}
\end{lstlisting}

目录条目由 \verb|\addcontentsline| 机制收集（见第五章
\texttt{.toc} 格式），样式定制只需改上述命令。

\section{习题与答案机制}

模板提供习题与答案配套的写作方式：

\begin{lstlisting}
\begin{exercise}[单调有界]
证明单调有界数列必收敛。
\end{exercise}

\begin{solution}
由确界原理：设上界存在，则上确界即极限。
\end{solution}
\end{lstlisting}

若想发布\emph{不含答案}的版本，只需在文档类选项中加入
\texttt{result=noanswer}：
\begin{lstlisting}
\documentclass[lang=cn, result=noanswer]{elegantbook}
\end{lstlisting}
此时 \texttt{solution} 与 \texttt{proof} 环境的内容会被自动
\emph{注释掉}（cls 内部用 \texttt{comment} 宏包的
\verb|\excludecomment| 实现），习题册与教师版可共用一份源码。
章末习题集使用 \texttt{problemset} 环境：

\begin{lstlisting}
\begin{problemset}
\item 习题一...
\item 习题二...
\end{problemset}
\end{lstlisting}

\section{双语机制}

\texttt{lang} 选项不仅切换 ctex/babel，还切换模板内部的
\textbf{名称宏}。例如中文模式下定义
\verb|\theoremname|、\verb|\definitionname| 等，英文模式下
分别定义为 Theorem、Definition。这些名称宏统一约定为
"\texttt{环境名 + name}"：

\begin{table}[htbp]
  \centering
  \caption{语言名称宏约定}\label{tab:names}
  \begin{tabular}{lll}
    \toprule
    名称宏 & 中文 & 英文 \\
    \midrule
    \texttt{\textbackslash theoremname} & 定理 & Theorem \\
    \texttt{\textbackslash definitionname} & 定义 & Definition \\
    \texttt{\textbackslash lemmname} & 引理 & Lemma \\
    \texttt{\textbackslash proofname} & 证明 & Proof \\
    \texttt{\textbackslash introductionname} & 本章导言 & Introduction \\
    \bottomrule
  \end{tabular}
\end{table}

这也是第四章 \verb|\elegantnewtheorem{...}| 要求环境名对应的
名称宏已定义的原因（未定义会触发文档类错误）。若使用模板未
覆盖的语言，可自行定义这些名称宏后配合 \texttt{lang=cn}
使用。

\section{分章编译与多人协作}

大书写作时推荐 \texttt{includeonly} 只编译当前章节，大幅
缩短单次编译时间：

\begin{lstlisting}
% 主文件导言区：只编译第二章
\includeonly{chapters/ch02_quickstart}
% ============================================================================
%  第一章  功能与模块划分
% ============================================================================

\chapter{功能与模块划分}

\section{什么是 ElegantBook}

ElegantBook 是一个基于 LaTeX 标准文档类 \texttt{book} 设计的优雅书籍模板，
由 ElegantLaTeX 团队维护，项目托管于
\href{https://github.com/ElegantLaTeX/ElegantBook}{GitHub}，
以 LaTeX Project Public License（LPPL）发布，可以自由使用与二次分发。
自 2013 年诞生以来，它经历了多次大版本重构，本文写作时对应的版本为
v4.6（\texttt{elegantbook.cls} 的文件头中通过 \verb|\ProvidesClass| 声明版本号）。

与 \texttt{book} 类相比，ElegantBook 的定位是\textbf{开箱即用的书籍排版方案}：
用户只需要填写标题、作者等元信息，选择喜欢的配色与版式选项，
即可得到包含封面、目录、章节、定理盒子、页眉页脚、参考文献在内的完整书籍，
而无需手动组装十几行宏包调用。

\begin{remark}
本手册约定：\textbf{不显式使用 ctex}。也就是说，文档源码中既不写
\verb|\usepackage{ctex}|，也不使用 \texttt{ctexbook}/\texttt{ctexart} 文档类；
中文排版所需的全部机制（标点压缩、中文字体、段落缩进等）由
ElegantBook 在 \texttt{lang=cn} 选项下于内部加载完成。
\end{remark}

\section{核心功能一览}

ElegantBook 的核心功能可以归纳为以下几类：

\begin{enumerate}
  \item \textbf{多语言支持}：通过 \texttt{lang} 选项一键切换
        中、英、意、法、荷、匈、德、蒙、葡、西、日等语言环境；
  \item \textbf{配色主题}：内置绿、青、蓝、灰、黑五套主题色，
        通过 \texttt{color} 选项选择，并可自行 \verb|\definecolor| 扩展；
  \item \textbf{封面自动生成}：由 \verb|\title|、\verb|\author|、
        \verb|\institute|、\verb|\cover| 等元信息自动排版出整页封面，
        无需手写封面环境；
  \item \textbf{定理盒子体系}：基于 \texttt{amsthm} 与 \texttt{tcolorbox}
        实现带彩色边框的定理、定义、引理、推论、命题、例题、证明等环境；
  \item \textbf{页眉页脚}：基于 \texttt{fancyhdr} 定制
        章节名、书名、页码的组合页眉，并支持书脊式页眉装饰；
  \item \textbf{图表排版}：内置 \texttt{graphicx}、\texttt{booktabs}、
        \texttt{caption}、\texttt{multicol}、\texttt{multirow} 等，
        图题与表题自动着色；
  \item \textbf{代码高亮}：内置 \texttt{listings} 样式，
        提供适合中文注释的等宽字体配置；
  \item \textbf{参考文献}：内置 \texttt{biblatex}，
        默认数值引用风格（\texttt{numeric-comp}），配合 \texttt{biber}
        或 \texttt{bibtex} 后端；
  \item \textbf{超链接与书签}：内置 \texttt{hyperref}，
        目录、交叉引用、外部链接全部可点击跳转；
  \item \textbf{三套版式}：\texttt{mode=fancy} 华丽版式（默认）、
        \texttt{mode=simple} 简约版式；\texttt{device=pad} 可为平板
        阅读优化版面。
\end{enumerate}

\section{模块划分}

从实现角度看，\texttt{elegantbook.cls} 是一个约一千五百行的文档类，
其内部可以划分为若干逻辑模块，如图~\ref{fig:modules} 所示。

\begin{figure}[htbp]
  \centering
  \begin{tikzpicture}[font=\small, >=stealth,
      lay/.style={draw=structurecolor, thick, rounded corners,
                  align=center, minimum width=11.5cm, fill=structurecolor!8},
      mod/.style={draw=main, thick, rounded corners, align=center,
                  minimum width=3.4cm, minimum height=1.1cm, fill=main!8}]
    % 顶层：用户接口
    \node[lay] (user) at (0,4.2) {用户接口：\texttt{\string\documentclass[\,\ldots\,]\{elegantbook\}}\\[2pt]
      {\scriptsize main.tex \quad chapters/ \quad refs.bib \quad figures/}};
    % 中间层：功能模块
    \node[mod] (lang)  at (-4.2, 1.6) {语言与字体\\ {\scriptsize ctex, fontspec}};
    \node[mod] (layo)  at (-1.4, 1.6) {版面样式\\ {\scriptsize geometry, fancyhdr, titlesec}};
    \node[mod] (math)  at (1.4, 1.6)  {数学与定理\\ {\scriptsize amsmath, amsthm, tcolorbox}};
    \node[mod] (graph) at (4.2, 1.6)  {图形与图表\\ {\scriptsize tikz, graphicx, booktabs}};
    \node[mod] (code)  at (-2.8,-1.2) {代码与列表\\ {\scriptsize listings, enumitem}};
    \node[mod] (bib)   at (0,  -1.2)  {参考文献\\ {\scriptsize biblatex}};
    \node[mod] (link)  at (2.8, -1.2) {超链接\\ {\scriptsize hyperref}};
    % 底层：文档类基座
    \node[lay] (base) at (0,-3.4) {基础层：\texttt{book} 类 \quad·\quad
      {\scriptsize xcolor, kvoptions, etoolbox, calc, setspace, csquotes 等}};
    % 连线
    \draw[->, thick, structurecolor] (user.south) -- (0, 2.6);
    \draw[->, thick, main] (lang.south) -- (-2.8,-1.9);
    \draw[->, thick, main] (layo.south) -- (-1.4,-1.9);
    \draw[->, thick, main] (math.south) -- (0,-1.9);
    \draw[->, thick, main] (graph.south) -- (2.8,-1.9);
    \draw[->, thick, second] (code.south) -- (-1.4,-2.7);
    \draw[->, thick, second] (bib.south) -- (0,-2.7);
    \draw[->, thick, second] (link.south) -- (1.4,-2.7);
  \end{tikzpicture}
  \caption{ElegantBook 的逻辑模块划分}\label{fig:modules}
\end{figure}

各模块的职责简述如下：

\begin{description}
  \item[语言与字体模块] 负责中文（\texttt{lang=cn} 时加载 ctex）、西文
        与日文（\texttt{luatexja}）排版；在 XeLaTeX 下通过
        \texttt{fontspec} 配置衬线、无衬线与等宽字体；
  \item[版面样式模块] 通过 \texttt{geometry} 设置页面边距，
        \texttt{titlesec} 定制章/节标题样式，\texttt{fancyhdr}
        绘制页眉页脚，\texttt{tocloft} 控制目录排版；
  \item[数学与定理模块] 加载 \texttt{amsmath} 系列宏包，
        用 \texttt{amsthm} 定义计数器，再用 \texttt{tcolorbox}
        把每个定理环境包装成彩色圆角盒子（详见第四章）；
  \item[图形与图表模块] 加载 \texttt{graphicx}、\texttt{tikz}、
        \texttt{booktabs}、\texttt{multicol}、\texttt{multirow} 等，
        并统一 \texttt{caption} 标签样式；
  \item[代码与列表模块] 提供 \texttt{lstlisting} 代码环境，
        \texttt{enumerate}/\texttt{enumitem} 定制化列表；
  \item[参考文献模块] 按 \texttt{citestyle}/\texttt{bibstyle} 选项
        配置 biblatex，支持 biber/bibtex 两种后端；
  \item[超链接模块] 统一配置 hyperref，保证书签、目录与引用可跳转。
\end{description}

\section{一个典型项目的文件组成}

一个基于 ElegantBook 的书籍项目通常由以下文件组成：

\begin{table}[htbp]
  \centering
  \caption{ElegantBook 项目的典型文件结构}\label{tab:project}
  \begin{tabular}{lll}
    \toprule
    文件 & 作用 & 是否必须 \\
    \midrule
    \texttt{main.tex}     & 主文件：文档类、封面元信息、\texttt{\string\include} 导入 & 必须 \\
    \texttt{chapters/}    & 章节目录，一章一文件（\texttt{ch01\_xxx.tex} 等） & 必须 \\
    \texttt{refs.bib}     & Bib\TeX 参考文献数据库 & 可选 \\
    \texttt{figures/}     & 插图目录（PNG/PDF/TikZ 源文件） & 可选 \\
    \texttt{elegantbook.cls} & 文档类文件 & 随发行版安装即可 \\
    \texttt{main.aux/.toc/.bbl/.out} & 编译中间文件（自动生成） & 否 \\
    \bottomrule
  \end{tabular}
\end{table}

编译中间文件由引擎自动生成，通常在最终交付时删除，需要时用
\texttt{latexmk -c} 或手动清理。

\section{本手册的结构}

本手册按照"先总览、再入门、后原理"的顺序组织：

\begin{itemize}
  \item 第一章（本章）：介绍功能与模块划分；
  \item 第二章：从最小示例出发，给出完整的入门流程；
  \item 第三章：版面与排版的数学原理与公式推导；
  \item 第四章：定理环境的内部设计与算法实现；
  \item 第五章：编译管线与文档类的内部自动化机制；
  \item 第六章：进阶定制方法；
  \item 附录 A：文档类选项速查表。
\end{itemize}

阅读完前两章即可上手使用；若希望深入理解模板的工作原理，甚至为其
贡献代码，请重点阅读第三、四、五章。

\section{历史与版本演进}

ElegantBook 自 2013 年发布以来经历了多次大版本重构，
理解其演进脉络有助于读懂网上的旧教程：

\begin{table}[htbp]
  \centering
  \caption{ElegantBook 主要版本演进}\label{tab:history}
  \begin{tabular}{llp{7.6cm}}
    \toprule
    版本 & 时间 & 主要变化 \\
    \midrule
    v1.x & 2013--2018 & 面向理工科笔记的简单书籍模板，
          提供封面与基础定理环境 \\
    v2   & 2018      & 引入主题色（绿/青/蓝）与
          \texttt{fancy} 版式，封面可插图 \\
    v3   & 2020      & 用 tcolorbox 重构全部定理环境，
          支持跨页断行的彩色盒子 \\
    v4   & 2021 至今 & 选项键值化（kvoptions）、多语言
          （cn/en/fr/de/jp 等）、\texttt{simple} 简约模式、
          \texttt{device=pad} 平板版式、用户自定义定理环境
          \texttt{\textbackslash elegantnewtheorem} \\
    \bottomrule
  \end{tabular}
\end{table}

当前主线版本（本文对应 v4.6）在设计上遵循三条原则：
\textbf{选项即配置}（一切外观差异都由文档类选项表达）、
\textbf{用户零侵入}（封面、目录、页眉等自动化，
用户只需填元信息与正文）、\textbf{可扩展}（颜色、定理环境、
字体均可覆盖）。

\section{特性的代码示例}

下面把第一节的功能清单落到具体代码上，帮助读者建立直观印象。

\subsection{封面元信息}

\begin{lstlisting}
\title{数学分析讲义}
\subtitle{微积分与级数理论}
\author{张三 \and 李四}
\institute{某某大学数学学院}
\date{2026 年 8 月}
\version{1.0}
\extrainfo{仅供课堂使用}
\logo{figures/logo.pdf}     % 可选：封面左上角校徽
\cover{figures/cover.pdf}   % 可选：整页封面背景图
\end{lstlisting}

\verb|\author| 支持 \verb|\and| 分隔多个作者；
\verb|\cover| 若给出，将作为整页背景铺满封面。

\subsection{主题色}

\begin{lstlisting}
\documentclass[lang=cn, color=cyan]{elegantbook}   % 青色调
\documentclass[lang=cn, color=green]{elegantbook}  % 绿色调
\end{lstlisting}

也可以完全自定义（见第六章）。

\subsection{定理盒子}

\begin{lstlisting}
\begin{theorem}[欧拉公式]
对任意实数 $x$，有 $e^{ix}=\cos x+i\sin x$。
\end{theorem}

\begin{lemma}
设 $f$ 在闭区间 $[a,b]$ 上连续，则 $f$ 一致连续。
\end{lemma}
\end{lstlisting}

\subsection{参考文献}

\begin{lstlisting}
% 主文件导言区
\addbibresource{refs.bib}
% 正文中引用
\cite{knuth1984texbook}
% 文末输出
\printbibliography[heading=bibintoc, title={参考文献}]
\end{lstlisting}

\section{模块内部宏包明细}

第一节给出了宏观模块划分，下表列出每个模块实际加载的核心宏包
及其职责，便于读者判断"模板内部已经替我做了什么"：

\begin{table}[htbp]
  \centering
  \caption{模块与内部宏包明细}\label{tab:packages}
  \begin{tabular}{llp{6.5cm}}
    \toprule
    模块 & 宏包 & 职责 \\
    \midrule
    语言 & ctex & 中文标点、缩进、字体配置（\texttt{lang=cn} 时）\\
    \multirow{2}{*}{字体} & fontspec & XeLaTeX 下系统字体加载 \\
                          & newtxtext/newtxmath & Times 风格西文与数学字体 \\
    \multirow{2}{*}{版面} & geometry & 页边距（见第三章推导）\\
                          & titlesec, fancyhdr, tocloft & 标题、页眉页脚、目录样式 \\
    \multirow{2}{*}{数学} & amsmath 系列 & 公式环境、符号字体 \\
                          & bm, esint & 粗斜体、积分符号 \\
    \multirow{2}{*}{盒子} & tcolorbox & 定理盒子的绘制引擎 \\
                          & amsthm & 定理计数器与 proof 环境 \\
    图形图表 & tikz, graphicx, booktabs, multicol, multirow & 矢量图、插图、三线表 \\
    代码列表 & listings, fancyvrb, enumerate, enumitem & 代码高亮与列表 \\
    文献 & biblatex & 参考文献（后端 biber/bibtex）\\
    工具 & hyperref, xcolor, kvoptions, etoolbox & 链接、颜色、选项解析 \\
    \bottomrule
  \end{tabular}
\end{table}

\begin{remark}
表中大部分宏包读者无需自行加载；若需要\emph{覆盖}其配置，
应在主文件导言区于 \verb|\begin{document}| 之前重新设置
（例如用 \verb|\definecolor| 覆盖主题色）。
\end{remark}

\section{与其他排版方案的对比}

选型时常见的问题："我为什么不用 ctexbook 或直接堆宏包？"
下表给出粗粒度对比：

\begin{table}[htbp]
  \centering
  \caption{常见中文书籍排版方案对比}\label{tab:compare}
  \begin{tabular}{lcccc}
    \toprule
    方案 & 中文开箱 & 彩色定理盒子 & 自动封面 & 配置成本 \\
    \midrule
    \texttt{ctexbook} + 手写宏包 & 好 & 无，需自建 & 无，需自建 & 高 \\
    \texttt{tufte-book} & 一般 & 无 & 无 & 中 \\
    \texttt{memoir} & 一般 & 无 & 无 & 高 \\
    \textbf{ElegantBook} & 好 & 内置 & 内置 & 低 \\
    \bottomrule
  \end{tabular}
\end{table}

ElegantBook 的取舍是：牺牲一部分"从零控制的自由度"，换取
"两天出成品"的确定性。对需要深度定制的用户，模板同样开放了
颜色、字体、定理环境与页眉页脚的覆盖接口（见第六、七章
与附录~A）。

\section{运行环境要求}

\begin{itemize}
  \item \textbf{引擎}：XeLaTeX（中文必需）；纯英文文档亦可用
        pdfLaTeX；
  \item \textbf{发行版}：TeX~Live（推荐，2021 版以后）或 MiKTeX；
        安装 \texttt{elegantbook} 宏包：
\begin{lstlisting}
tlmgr install elegantbook
\end{lstlisting}
  \item \textbf{参考文献后端}：biber（默认）或 BibTeX；
  \item \textbf{字体}：中文排版需要系统中存在中文字体
        （Windows 自带宋体/黑体/楷体，macOS 自带苹方/宋体等）；
  \item \textbf{编辑工具}：TeXstudio、VS Code + LaTeX~Workshop、
        Overleaf（在线）均可；编辑器内配置"XeLaTeX 编译 +
        biber 参考文献"即可（见第二章）。
\end{itemize}

% ============================================================================
%  第二章  快速入门
% ============================================================================

\chapter{快速入门}

本章从一个最小示例出发，介绍从零开始编写一份 ElegantBook 中文文档的
完整流程：源文件怎么写、怎么编译、常用环境怎么用、出现问题怎么排错。
读者只需跟随本章操作即可产出第一份成品。

\section{一个最小示例}

下面是最小的完整源文件（假设文件名为 \texttt{main.tex}）：

\begin{lstlisting}
% main.tex —— 最小示例
\documentclass[lang=cn]{elegantbook}

\title{我的第一本书}
\author{张三}
\institute{某某大学}
\date{\today}

\begin{document}
\maketitle
\tableofcontents

\chapter{引言}
这是一本用 ElegantBook 排版的书籍。
\end{document}
\end{lstlisting}

将以上内容保存为 \texttt{main.tex} 后，依次执行第~\ref{sec:compile}~节给出的
四条命令，即可得到 \texttt{main.pdf}：其中封面、目录、章节标题的排版
全部由文档类自动完成。可以看到，用户只需要填写元信息与正文内容，
其余工作（封面绘制、页眉页脚、目录生成、超链接着色等）都由模板代劳。

\section{编译流程与引擎要求}\label{sec:compile}

ElegantBook 的中文模式需要 \textbf{XeLaTeX} 引擎（它原生支持 UTF-8 与系统
字体），因此请使用 XeLaTeX 编译；若使用 \texttt{lang=en} 的纯英文文档，
则 pdfLaTeX 亦可。推荐的完整编译命令为：

\begin{verbatim}
xelatex main.tex      % 第 1 遍：生成 .aux、.toc、.bcf 等辅助文件
biber   main          % 处理参考文献：读 .bcf，写 .bbl
xelatex main.tex      % 第 2 遍：引用参考文献并写入目录
xelatex main.tex      % 第 3 遍：收敛页码与交叉引用
\end{verbatim}

为什么需要多遍编译？因为 LaTeX 采用"两遍/多遍编译"模型：第一遍只知道
每个\emph{标签}（\verb|\label|）的位置，但页码与编号要等排版完成才能确定；
目录（\texttt{.toc}）、交叉引用（\texttt{.aux}）和参考文献（\texttt{.bbl}）
正是在各遍之间传递信息的文件。若正文或引用有变动，第三遍通常已足够收敛。

\begin{remark}
也可以使用 \texttt{latexmk} 自动完成上述流程：
\begin{verbatim}
latexmk -xelatex -interaction=nonstopmode main.tex
\end{verbatim}
该工具会监测源文件变化并自动重跑所需遍数，适合写作阶段使用。
\end{remark}

\section{文档类选项}

\verb|\documentclass| 的可选参数以键值对形式给出，最常用的组合如下：

\begin{lstlisting}
\documentclass[11pt,a4paper,lang=cn,color=blue,mode=fancy,
               device=normal,math=cm,titlestyle=hang]{elegantbook}
\end{lstlisting}

其中各选项的含义为：

\begin{description}
  \item[\texttt{lang=cn}] 语言环境，中文排版；可选 \texttt{en}（英文）、
        \texttt{fr}（法文）、\texttt{de}（德文）、\texttt{jp}（日文）等；
  \item[\texttt{color=blue}] 主题色，可选 \texttt{green}、\texttt{cyan}、
        \texttt{blue}、\texttt{gray}、\texttt{black}、\texttt{nocolor}；
  \item[\texttt{mode=fancy}] 版式，\texttt{fancy} 为华丽版（默认），
        \texttt{simple} 为简约版；
  \item[\texttt{device=normal}] 输出设备，\texttt{normal}（默认）或
        \texttt{pad}（平板优化）；
  \item[\texttt{math=cm}] 数学字体，\texttt{cm}（默认，Computer Modern）、
        \texttt{newtx}（Times 风格）或 \texttt{mtpro2}；
  \item[\texttt{titlestyle=hang}] 章节标题样式，\texttt{hang}（悬挂式，
        默认）或 \texttt{display}（通栏式）。
\end{description}

其余选项（如 \texttt{chinesefont}、\texttt{toc}、\texttt{thmcnt}、
\texttt{citestyle}、\texttt{bibend} 等）在附录~A 中给出完整速查表。

\section{封面与前置部分}

在正文之前，需要填写封面元信息。ElegantBook 提供的封面命令包括：

\begin{lstlisting}
\title{书名}                 % 必须
\subtitle{副标题}            % 可选
\author{作者}                % 必须
\institute{单位/学校}        % 可选
\date{2026 年 8 月}          % 可选，默认当天
\version{1.0}                % 可选，显示在封面
\extrainfo{副信息}           % 可选
\logo{figures/logo.pdf}      % 可选，封面校徽
\cover{figures/cover.pdf}    % 可选，整页封面图
\keywords{关键词}            % 可选
\end{lstlisting}

然后在 \verb|\begin{document}| 之后依次调用
\verb|\maketitle|、\verb|\tableofcontents|，并用
\verb|\frontmatter|/\verb|\mainmatter| 划分前置部分与正文：

\begin{lstlisting}
\begin{document}
\maketitle
\frontmatter
\tableofcontents
\mainmatter
% ============================================================================
%  第一章  功能与模块划分
% ============================================================================

\chapter{功能与模块划分}

\section{什么是 ElegantBook}

ElegantBook 是一个基于 LaTeX 标准文档类 \texttt{book} 设计的优雅书籍模板，
由 ElegantLaTeX 团队维护，项目托管于
\href{https://github.com/ElegantLaTeX/ElegantBook}{GitHub}，
以 LaTeX Project Public License（LPPL）发布，可以自由使用与二次分发。
自 2013 年诞生以来，它经历了多次大版本重构，本文写作时对应的版本为
v4.6（\texttt{elegantbook.cls} 的文件头中通过 \verb|\ProvidesClass| 声明版本号）。

与 \texttt{book} 类相比，ElegantBook 的定位是\textbf{开箱即用的书籍排版方案}：
用户只需要填写标题、作者等元信息，选择喜欢的配色与版式选项，
即可得到包含封面、目录、章节、定理盒子、页眉页脚、参考文献在内的完整书籍，
而无需手动组装十几行宏包调用。

\begin{remark}
本手册约定：\textbf{不显式使用 ctex}。也就是说，文档源码中既不写
\verb|\usepackage{ctex}|，也不使用 \texttt{ctexbook}/\texttt{ctexart} 文档类；
中文排版所需的全部机制（标点压缩、中文字体、段落缩进等）由
ElegantBook 在 \texttt{lang=cn} 选项下于内部加载完成。
\end{remark}

\section{核心功能一览}

ElegantBook 的核心功能可以归纳为以下几类：

\begin{enumerate}
  \item \textbf{多语言支持}：通过 \texttt{lang} 选项一键切换
        中、英、意、法、荷、匈、德、蒙、葡、西、日等语言环境；
  \item \textbf{配色主题}：内置绿、青、蓝、灰、黑五套主题色，
        通过 \texttt{color} 选项选择，并可自行 \verb|\definecolor| 扩展；
  \item \textbf{封面自动生成}：由 \verb|\title|、\verb|\author|、
        \verb|\institute|、\verb|\cover| 等元信息自动排版出整页封面，
        无需手写封面环境；
  \item \textbf{定理盒子体系}：基于 \texttt{amsthm} 与 \texttt{tcolorbox}
        实现带彩色边框的定理、定义、引理、推论、命题、例题、证明等环境；
  \item \textbf{页眉页脚}：基于 \texttt{fancyhdr} 定制
        章节名、书名、页码的组合页眉，并支持书脊式页眉装饰；
  \item \textbf{图表排版}：内置 \texttt{graphicx}、\texttt{booktabs}、
        \texttt{caption}、\texttt{multicol}、\texttt{multirow} 等，
        图题与表题自动着色；
  \item \textbf{代码高亮}：内置 \texttt{listings} 样式，
        提供适合中文注释的等宽字体配置；
  \item \textbf{参考文献}：内置 \texttt{biblatex}，
        默认数值引用风格（\texttt{numeric-comp}），配合 \texttt{biber}
        或 \texttt{bibtex} 后端；
  \item \textbf{超链接与书签}：内置 \texttt{hyperref}，
        目录、交叉引用、外部链接全部可点击跳转；
  \item \textbf{三套版式}：\texttt{mode=fancy} 华丽版式（默认）、
        \texttt{mode=simple} 简约版式；\texttt{device=pad} 可为平板
        阅读优化版面。
\end{enumerate}

\section{模块划分}

从实现角度看，\texttt{elegantbook.cls} 是一个约一千五百行的文档类，
其内部可以划分为若干逻辑模块，如图~\ref{fig:modules} 所示。

\begin{figure}[htbp]
  \centering
  \begin{tikzpicture}[font=\small, >=stealth,
      lay/.style={draw=structurecolor, thick, rounded corners,
                  align=center, minimum width=11.5cm, fill=structurecolor!8},
      mod/.style={draw=main, thick, rounded corners, align=center,
                  minimum width=3.4cm, minimum height=1.1cm, fill=main!8}]
    % 顶层：用户接口
    \node[lay] (user) at (0,4.2) {用户接口：\texttt{\string\documentclass[\,\ldots\,]\{elegantbook\}}\\[2pt]
      {\scriptsize main.tex \quad chapters/ \quad refs.bib \quad figures/}};
    % 中间层：功能模块
    \node[mod] (lang)  at (-4.2, 1.6) {语言与字体\\ {\scriptsize ctex, fontspec}};
    \node[mod] (layo)  at (-1.4, 1.6) {版面样式\\ {\scriptsize geometry, fancyhdr, titlesec}};
    \node[mod] (math)  at (1.4, 1.6)  {数学与定理\\ {\scriptsize amsmath, amsthm, tcolorbox}};
    \node[mod] (graph) at (4.2, 1.6)  {图形与图表\\ {\scriptsize tikz, graphicx, booktabs}};
    \node[mod] (code)  at (-2.8,-1.2) {代码与列表\\ {\scriptsize listings, enumitem}};
    \node[mod] (bib)   at (0,  -1.2)  {参考文献\\ {\scriptsize biblatex}};
    \node[mod] (link)  at (2.8, -1.2) {超链接\\ {\scriptsize hyperref}};
    % 底层：文档类基座
    \node[lay] (base) at (0,-3.4) {基础层：\texttt{book} 类 \quad·\quad
      {\scriptsize xcolor, kvoptions, etoolbox, calc, setspace, csquotes 等}};
    % 连线
    \draw[->, thick, structurecolor] (user.south) -- (0, 2.6);
    \draw[->, thick, main] (lang.south) -- (-2.8,-1.9);
    \draw[->, thick, main] (layo.south) -- (-1.4,-1.9);
    \draw[->, thick, main] (math.south) -- (0,-1.9);
    \draw[->, thick, main] (graph.south) -- (2.8,-1.9);
    \draw[->, thick, second] (code.south) -- (-1.4,-2.7);
    \draw[->, thick, second] (bib.south) -- (0,-2.7);
    \draw[->, thick, second] (link.south) -- (1.4,-2.7);
  \end{tikzpicture}
  \caption{ElegantBook 的逻辑模块划分}\label{fig:modules}
\end{figure}

各模块的职责简述如下：

\begin{description}
  \item[语言与字体模块] 负责中文（\texttt{lang=cn} 时加载 ctex）、西文
        与日文（\texttt{luatexja}）排版；在 XeLaTeX 下通过
        \texttt{fontspec} 配置衬线、无衬线与等宽字体；
  \item[版面样式模块] 通过 \texttt{geometry} 设置页面边距，
        \texttt{titlesec} 定制章/节标题样式，\texttt{fancyhdr}
        绘制页眉页脚，\texttt{tocloft} 控制目录排版；
  \item[数学与定理模块] 加载 \texttt{amsmath} 系列宏包，
        用 \texttt{amsthm} 定义计数器，再用 \texttt{tcolorbox}
        把每个定理环境包装成彩色圆角盒子（详见第四章）；
  \item[图形与图表模块] 加载 \texttt{graphicx}、\texttt{tikz}、
        \texttt{booktabs}、\texttt{multicol}、\texttt{multirow} 等，
        并统一 \texttt{caption} 标签样式；
  \item[代码与列表模块] 提供 \texttt{lstlisting} 代码环境，
        \texttt{enumerate}/\texttt{enumitem} 定制化列表；
  \item[参考文献模块] 按 \texttt{citestyle}/\texttt{bibstyle} 选项
        配置 biblatex，支持 biber/bibtex 两种后端；
  \item[超链接模块] 统一配置 hyperref，保证书签、目录与引用可跳转。
\end{description}

\section{一个典型项目的文件组成}

一个基于 ElegantBook 的书籍项目通常由以下文件组成：

\begin{table}[htbp]
  \centering
  \caption{ElegantBook 项目的典型文件结构}\label{tab:project}
  \begin{tabular}{lll}
    \toprule
    文件 & 作用 & 是否必须 \\
    \midrule
    \texttt{main.tex}     & 主文件：文档类、封面元信息、\texttt{\string\include} 导入 & 必须 \\
    \texttt{chapters/}    & 章节目录，一章一文件（\texttt{ch01\_xxx.tex} 等） & 必须 \\
    \texttt{refs.bib}     & Bib\TeX 参考文献数据库 & 可选 \\
    \texttt{figures/}     & 插图目录（PNG/PDF/TikZ 源文件） & 可选 \\
    \texttt{elegantbook.cls} & 文档类文件 & 随发行版安装即可 \\
    \texttt{main.aux/.toc/.bbl/.out} & 编译中间文件（自动生成） & 否 \\
    \bottomrule
  \end{tabular}
\end{table}

编译中间文件由引擎自动生成，通常在最终交付时删除，需要时用
\texttt{latexmk -c} 或手动清理。

\section{本手册的结构}

本手册按照"先总览、再入门、后原理"的顺序组织：

\begin{itemize}
  \item 第一章（本章）：介绍功能与模块划分；
  \item 第二章：从最小示例出发，给出完整的入门流程；
  \item 第三章：版面与排版的数学原理与公式推导；
  \item 第四章：定理环境的内部设计与算法实现；
  \item 第五章：编译管线与文档类的内部自动化机制；
  \item 第六章：进阶定制方法；
  \item 附录 A：文档类选项速查表。
\end{itemize}

阅读完前两章即可上手使用；若希望深入理解模板的工作原理，甚至为其
贡献代码，请重点阅读第三、四、五章。

\section{历史与版本演进}

ElegantBook 自 2013 年发布以来经历了多次大版本重构，
理解其演进脉络有助于读懂网上的旧教程：

\begin{table}[htbp]
  \centering
  \caption{ElegantBook 主要版本演进}\label{tab:history}
  \begin{tabular}{llp{7.6cm}}
    \toprule
    版本 & 时间 & 主要变化 \\
    \midrule
    v1.x & 2013--2018 & 面向理工科笔记的简单书籍模板，
          提供封面与基础定理环境 \\
    v2   & 2018      & 引入主题色（绿/青/蓝）与
          \texttt{fancy} 版式，封面可插图 \\
    v3   & 2020      & 用 tcolorbox 重构全部定理环境，
          支持跨页断行的彩色盒子 \\
    v4   & 2021 至今 & 选项键值化（kvoptions）、多语言
          （cn/en/fr/de/jp 等）、\texttt{simple} 简约模式、
          \texttt{device=pad} 平板版式、用户自定义定理环境
          \texttt{\textbackslash elegantnewtheorem} \\
    \bottomrule
  \end{tabular}
\end{table}

当前主线版本（本文对应 v4.6）在设计上遵循三条原则：
\textbf{选项即配置}（一切外观差异都由文档类选项表达）、
\textbf{用户零侵入}（封面、目录、页眉等自动化，
用户只需填元信息与正文）、\textbf{可扩展}（颜色、定理环境、
字体均可覆盖）。

\section{特性的代码示例}

下面把第一节的功能清单落到具体代码上，帮助读者建立直观印象。

\subsection{封面元信息}

\begin{lstlisting}
\title{数学分析讲义}
\subtitle{微积分与级数理论}
\author{张三 \and 李四}
\institute{某某大学数学学院}
\date{2026 年 8 月}
\version{1.0}
\extrainfo{仅供课堂使用}
\logo{figures/logo.pdf}     % 可选：封面左上角校徽
\cover{figures/cover.pdf}   % 可选：整页封面背景图
\end{lstlisting}

\verb|\author| 支持 \verb|\and| 分隔多个作者；
\verb|\cover| 若给出，将作为整页背景铺满封面。

\subsection{主题色}

\begin{lstlisting}
\documentclass[lang=cn, color=cyan]{elegantbook}   % 青色调
\documentclass[lang=cn, color=green]{elegantbook}  % 绿色调
\end{lstlisting}

也可以完全自定义（见第六章）。

\subsection{定理盒子}

\begin{lstlisting}
\begin{theorem}[欧拉公式]
对任意实数 $x$，有 $e^{ix}=\cos x+i\sin x$。
\end{theorem}

\begin{lemma}
设 $f$ 在闭区间 $[a,b]$ 上连续，则 $f$ 一致连续。
\end{lemma}
\end{lstlisting}

\subsection{参考文献}

\begin{lstlisting}
% 主文件导言区
\addbibresource{refs.bib}
% 正文中引用
\cite{knuth1984texbook}
% 文末输出
\printbibliography[heading=bibintoc, title={参考文献}]
\end{lstlisting}

\section{模块内部宏包明细}

第一节给出了宏观模块划分，下表列出每个模块实际加载的核心宏包
及其职责，便于读者判断"模板内部已经替我做了什么"：

\begin{table}[htbp]
  \centering
  \caption{模块与内部宏包明细}\label{tab:packages}
  \begin{tabular}{llp{6.5cm}}
    \toprule
    模块 & 宏包 & 职责 \\
    \midrule
    语言 & ctex & 中文标点、缩进、字体配置（\texttt{lang=cn} 时）\\
    \multirow{2}{*}{字体} & fontspec & XeLaTeX 下系统字体加载 \\
                          & newtxtext/newtxmath & Times 风格西文与数学字体 \\
    \multirow{2}{*}{版面} & geometry & 页边距（见第三章推导）\\
                          & titlesec, fancyhdr, tocloft & 标题、页眉页脚、目录样式 \\
    \multirow{2}{*}{数学} & amsmath 系列 & 公式环境、符号字体 \\
                          & bm, esint & 粗斜体、积分符号 \\
    \multirow{2}{*}{盒子} & tcolorbox & 定理盒子的绘制引擎 \\
                          & amsthm & 定理计数器与 proof 环境 \\
    图形图表 & tikz, graphicx, booktabs, multicol, multirow & 矢量图、插图、三线表 \\
    代码列表 & listings, fancyvrb, enumerate, enumitem & 代码高亮与列表 \\
    文献 & biblatex & 参考文献（后端 biber/bibtex）\\
    工具 & hyperref, xcolor, kvoptions, etoolbox & 链接、颜色、选项解析 \\
    \bottomrule
  \end{tabular}
\end{table}

\begin{remark}
表中大部分宏包读者无需自行加载；若需要\emph{覆盖}其配置，
应在主文件导言区于 \verb|\begin{document}| 之前重新设置
（例如用 \verb|\definecolor| 覆盖主题色）。
\end{remark}

\section{与其他排版方案的对比}

选型时常见的问题："我为什么不用 ctexbook 或直接堆宏包？"
下表给出粗粒度对比：

\begin{table}[htbp]
  \centering
  \caption{常见中文书籍排版方案对比}\label{tab:compare}
  \begin{tabular}{lcccc}
    \toprule
    方案 & 中文开箱 & 彩色定理盒子 & 自动封面 & 配置成本 \\
    \midrule
    \texttt{ctexbook} + 手写宏包 & 好 & 无，需自建 & 无，需自建 & 高 \\
    \texttt{tufte-book} & 一般 & 无 & 无 & 中 \\
    \texttt{memoir} & 一般 & 无 & 无 & 高 \\
    \textbf{ElegantBook} & 好 & 内置 & 内置 & 低 \\
    \bottomrule
  \end{tabular}
\end{table}

ElegantBook 的取舍是：牺牲一部分"从零控制的自由度"，换取
"两天出成品"的确定性。对需要深度定制的用户，模板同样开放了
颜色、字体、定理环境与页眉页脚的覆盖接口（见第六、七章
与附录~A）。

\section{运行环境要求}

\begin{itemize}
  \item \textbf{引擎}：XeLaTeX（中文必需）；纯英文文档亦可用
        pdfLaTeX；
  \item \textbf{发行版}：TeX~Live（推荐，2021 版以后）或 MiKTeX；
        安装 \texttt{elegantbook} 宏包：
\begin{lstlisting}
tlmgr install elegantbook
\end{lstlisting}
  \item \textbf{参考文献后端}：biber（默认）或 BibTeX；
  \item \textbf{字体}：中文排版需要系统中存在中文字体
        （Windows 自带宋体/黑体/楷体，macOS 自带苹方/宋体等）；
  \item \textbf{编辑工具}：TeXstudio、VS Code + LaTeX~Workshop、
        Overleaf（在线）均可；编辑器内配置"XeLaTeX 编译 +
        biber 参考文献"即可（见第二章）。
\end{itemize}

% ============================================================================
%  第二章  快速入门
% ============================================================================

\chapter{快速入门}

本章从一个最小示例出发，介绍从零开始编写一份 ElegantBook 中文文档的
完整流程：源文件怎么写、怎么编译、常用环境怎么用、出现问题怎么排错。
读者只需跟随本章操作即可产出第一份成品。

\section{一个最小示例}

下面是最小的完整源文件（假设文件名为 \texttt{main.tex}）：

\begin{lstlisting}
% main.tex —— 最小示例
\documentclass[lang=cn]{elegantbook}

\title{我的第一本书}
\author{张三}
\institute{某某大学}
\date{\today}

\begin{document}
\maketitle
\tableofcontents

\chapter{引言}
这是一本用 ElegantBook 排版的书籍。
\end{document}
\end{lstlisting}

将以上内容保存为 \texttt{main.tex} 后，依次执行第~\ref{sec:compile}~节给出的
四条命令，即可得到 \texttt{main.pdf}：其中封面、目录、章节标题的排版
全部由文档类自动完成。可以看到，用户只需要填写元信息与正文内容，
其余工作（封面绘制、页眉页脚、目录生成、超链接着色等）都由模板代劳。

\section{编译流程与引擎要求}\label{sec:compile}

ElegantBook 的中文模式需要 \textbf{XeLaTeX} 引擎（它原生支持 UTF-8 与系统
字体），因此请使用 XeLaTeX 编译；若使用 \texttt{lang=en} 的纯英文文档，
则 pdfLaTeX 亦可。推荐的完整编译命令为：

\begin{verbatim}
xelatex main.tex      % 第 1 遍：生成 .aux、.toc、.bcf 等辅助文件
biber   main          % 处理参考文献：读 .bcf，写 .bbl
xelatex main.tex      % 第 2 遍：引用参考文献并写入目录
xelatex main.tex      % 第 3 遍：收敛页码与交叉引用
\end{verbatim}

为什么需要多遍编译？因为 LaTeX 采用"两遍/多遍编译"模型：第一遍只知道
每个\emph{标签}（\verb|\label|）的位置，但页码与编号要等排版完成才能确定；
目录（\texttt{.toc}）、交叉引用（\texttt{.aux}）和参考文献（\texttt{.bbl}）
正是在各遍之间传递信息的文件。若正文或引用有变动，第三遍通常已足够收敛。

\begin{remark}
也可以使用 \texttt{latexmk} 自动完成上述流程：
\begin{verbatim}
latexmk -xelatex -interaction=nonstopmode main.tex
\end{verbatim}
该工具会监测源文件变化并自动重跑所需遍数，适合写作阶段使用。
\end{remark}

\section{文档类选项}

\verb|\documentclass| 的可选参数以键值对形式给出，最常用的组合如下：

\begin{lstlisting}
\documentclass[11pt,a4paper,lang=cn,color=blue,mode=fancy,
               device=normal,math=cm,titlestyle=hang]{elegantbook}
\end{lstlisting}

其中各选项的含义为：

\begin{description}
  \item[\texttt{lang=cn}] 语言环境，中文排版；可选 \texttt{en}（英文）、
        \texttt{fr}（法文）、\texttt{de}（德文）、\texttt{jp}（日文）等；
  \item[\texttt{color=blue}] 主题色，可选 \texttt{green}、\texttt{cyan}、
        \texttt{blue}、\texttt{gray}、\texttt{black}、\texttt{nocolor}；
  \item[\texttt{mode=fancy}] 版式，\texttt{fancy} 为华丽版（默认），
        \texttt{simple} 为简约版；
  \item[\texttt{device=normal}] 输出设备，\texttt{normal}（默认）或
        \texttt{pad}（平板优化）；
  \item[\texttt{math=cm}] 数学字体，\texttt{cm}（默认，Computer Modern）、
        \texttt{newtx}（Times 风格）或 \texttt{mtpro2}；
  \item[\texttt{titlestyle=hang}] 章节标题样式，\texttt{hang}（悬挂式，
        默认）或 \texttt{display}（通栏式）。
\end{description}

其余选项（如 \texttt{chinesefont}、\texttt{toc}、\texttt{thmcnt}、
\texttt{citestyle}、\texttt{bibend} 等）在附录~A 中给出完整速查表。

\section{封面与前置部分}

在正文之前，需要填写封面元信息。ElegantBook 提供的封面命令包括：

\begin{lstlisting}
\title{书名}                 % 必须
\subtitle{副标题}            % 可选
\author{作者}                % 必须
\institute{单位/学校}        % 可选
\date{2026 年 8 月}          % 可选，默认当天
\version{1.0}                % 可选，显示在封面
\extrainfo{副信息}           % 可选
\logo{figures/logo.pdf}      % 可选，封面校徽
\cover{figures/cover.pdf}    % 可选，整页封面图
\keywords{关键词}            % 可选
\end{lstlisting}

然后在 \verb|\begin{document}| 之后依次调用
\verb|\maketitle|、\verb|\tableofcontents|，并用
\verb|\frontmatter|/\verb|\mainmatter| 划分前置部分与正文：

\begin{lstlisting}
\begin{document}
\maketitle
\frontmatter
\tableofcontents
\mainmatter
% ============================================================================
%  第一章  功能与模块划分
% ============================================================================

\chapter{功能与模块划分}

\section{什么是 ElegantBook}

ElegantBook 是一个基于 LaTeX 标准文档类 \texttt{book} 设计的优雅书籍模板，
由 ElegantLaTeX 团队维护，项目托管于
\href{https://github.com/ElegantLaTeX/ElegantBook}{GitHub}，
以 LaTeX Project Public License（LPPL）发布，可以自由使用与二次分发。
自 2013 年诞生以来，它经历了多次大版本重构，本文写作时对应的版本为
v4.6（\texttt{elegantbook.cls} 的文件头中通过 \verb|\ProvidesClass| 声明版本号）。

与 \texttt{book} 类相比，ElegantBook 的定位是\textbf{开箱即用的书籍排版方案}：
用户只需要填写标题、作者等元信息，选择喜欢的配色与版式选项，
即可得到包含封面、目录、章节、定理盒子、页眉页脚、参考文献在内的完整书籍，
而无需手动组装十几行宏包调用。

\begin{remark}
本手册约定：\textbf{不显式使用 ctex}。也就是说，文档源码中既不写
\verb|\usepackage{ctex}|，也不使用 \texttt{ctexbook}/\texttt{ctexart} 文档类；
中文排版所需的全部机制（标点压缩、中文字体、段落缩进等）由
ElegantBook 在 \texttt{lang=cn} 选项下于内部加载完成。
\end{remark}

\section{核心功能一览}

ElegantBook 的核心功能可以归纳为以下几类：

\begin{enumerate}
  \item \textbf{多语言支持}：通过 \texttt{lang} 选项一键切换
        中、英、意、法、荷、匈、德、蒙、葡、西、日等语言环境；
  \item \textbf{配色主题}：内置绿、青、蓝、灰、黑五套主题色，
        通过 \texttt{color} 选项选择，并可自行 \verb|\definecolor| 扩展；
  \item \textbf{封面自动生成}：由 \verb|\title|、\verb|\author|、
        \verb|\institute|、\verb|\cover| 等元信息自动排版出整页封面，
        无需手写封面环境；
  \item \textbf{定理盒子体系}：基于 \texttt{amsthm} 与 \texttt{tcolorbox}
        实现带彩色边框的定理、定义、引理、推论、命题、例题、证明等环境；
  \item \textbf{页眉页脚}：基于 \texttt{fancyhdr} 定制
        章节名、书名、页码的组合页眉，并支持书脊式页眉装饰；
  \item \textbf{图表排版}：内置 \texttt{graphicx}、\texttt{booktabs}、
        \texttt{caption}、\texttt{multicol}、\texttt{multirow} 等，
        图题与表题自动着色；
  \item \textbf{代码高亮}：内置 \texttt{listings} 样式，
        提供适合中文注释的等宽字体配置；
  \item \textbf{参考文献}：内置 \texttt{biblatex}，
        默认数值引用风格（\texttt{numeric-comp}），配合 \texttt{biber}
        或 \texttt{bibtex} 后端；
  \item \textbf{超链接与书签}：内置 \texttt{hyperref}，
        目录、交叉引用、外部链接全部可点击跳转；
  \item \textbf{三套版式}：\texttt{mode=fancy} 华丽版式（默认）、
        \texttt{mode=simple} 简约版式；\texttt{device=pad} 可为平板
        阅读优化版面。
\end{enumerate}

\section{模块划分}

从实现角度看，\texttt{elegantbook.cls} 是一个约一千五百行的文档类，
其内部可以划分为若干逻辑模块，如图~\ref{fig:modules} 所示。

\begin{figure}[htbp]
  \centering
  \begin{tikzpicture}[font=\small, >=stealth,
      lay/.style={draw=structurecolor, thick, rounded corners,
                  align=center, minimum width=11.5cm, fill=structurecolor!8},
      mod/.style={draw=main, thick, rounded corners, align=center,
                  minimum width=3.4cm, minimum height=1.1cm, fill=main!8}]
    % 顶层：用户接口
    \node[lay] (user) at (0,4.2) {用户接口：\texttt{\string\documentclass[\,\ldots\,]\{elegantbook\}}\\[2pt]
      {\scriptsize main.tex \quad chapters/ \quad refs.bib \quad figures/}};
    % 中间层：功能模块
    \node[mod] (lang)  at (-4.2, 1.6) {语言与字体\\ {\scriptsize ctex, fontspec}};
    \node[mod] (layo)  at (-1.4, 1.6) {版面样式\\ {\scriptsize geometry, fancyhdr, titlesec}};
    \node[mod] (math)  at (1.4, 1.6)  {数学与定理\\ {\scriptsize amsmath, amsthm, tcolorbox}};
    \node[mod] (graph) at (4.2, 1.6)  {图形与图表\\ {\scriptsize tikz, graphicx, booktabs}};
    \node[mod] (code)  at (-2.8,-1.2) {代码与列表\\ {\scriptsize listings, enumitem}};
    \node[mod] (bib)   at (0,  -1.2)  {参考文献\\ {\scriptsize biblatex}};
    \node[mod] (link)  at (2.8, -1.2) {超链接\\ {\scriptsize hyperref}};
    % 底层：文档类基座
    \node[lay] (base) at (0,-3.4) {基础层：\texttt{book} 类 \quad·\quad
      {\scriptsize xcolor, kvoptions, etoolbox, calc, setspace, csquotes 等}};
    % 连线
    \draw[->, thick, structurecolor] (user.south) -- (0, 2.6);
    \draw[->, thick, main] (lang.south) -- (-2.8,-1.9);
    \draw[->, thick, main] (layo.south) -- (-1.4,-1.9);
    \draw[->, thick, main] (math.south) -- (0,-1.9);
    \draw[->, thick, main] (graph.south) -- (2.8,-1.9);
    \draw[->, thick, second] (code.south) -- (-1.4,-2.7);
    \draw[->, thick, second] (bib.south) -- (0,-2.7);
    \draw[->, thick, second] (link.south) -- (1.4,-2.7);
  \end{tikzpicture}
  \caption{ElegantBook 的逻辑模块划分}\label{fig:modules}
\end{figure}

各模块的职责简述如下：

\begin{description}
  \item[语言与字体模块] 负责中文（\texttt{lang=cn} 时加载 ctex）、西文
        与日文（\texttt{luatexja}）排版；在 XeLaTeX 下通过
        \texttt{fontspec} 配置衬线、无衬线与等宽字体；
  \item[版面样式模块] 通过 \texttt{geometry} 设置页面边距，
        \texttt{titlesec} 定制章/节标题样式，\texttt{fancyhdr}
        绘制页眉页脚，\texttt{tocloft} 控制目录排版；
  \item[数学与定理模块] 加载 \texttt{amsmath} 系列宏包，
        用 \texttt{amsthm} 定义计数器，再用 \texttt{tcolorbox}
        把每个定理环境包装成彩色圆角盒子（详见第四章）；
  \item[图形与图表模块] 加载 \texttt{graphicx}、\texttt{tikz}、
        \texttt{booktabs}、\texttt{multicol}、\texttt{multirow} 等，
        并统一 \texttt{caption} 标签样式；
  \item[代码与列表模块] 提供 \texttt{lstlisting} 代码环境，
        \texttt{enumerate}/\texttt{enumitem} 定制化列表；
  \item[参考文献模块] 按 \texttt{citestyle}/\texttt{bibstyle} 选项
        配置 biblatex，支持 biber/bibtex 两种后端；
  \item[超链接模块] 统一配置 hyperref，保证书签、目录与引用可跳转。
\end{description}

\section{一个典型项目的文件组成}

一个基于 ElegantBook 的书籍项目通常由以下文件组成：

\begin{table}[htbp]
  \centering
  \caption{ElegantBook 项目的典型文件结构}\label{tab:project}
  \begin{tabular}{lll}
    \toprule
    文件 & 作用 & 是否必须 \\
    \midrule
    \texttt{main.tex}     & 主文件：文档类、封面元信息、\texttt{\string\include} 导入 & 必须 \\
    \texttt{chapters/}    & 章节目录，一章一文件（\texttt{ch01\_xxx.tex} 等） & 必须 \\
    \texttt{refs.bib}     & Bib\TeX 参考文献数据库 & 可选 \\
    \texttt{figures/}     & 插图目录（PNG/PDF/TikZ 源文件） & 可选 \\
    \texttt{elegantbook.cls} & 文档类文件 & 随发行版安装即可 \\
    \texttt{main.aux/.toc/.bbl/.out} & 编译中间文件（自动生成） & 否 \\
    \bottomrule
  \end{tabular}
\end{table}

编译中间文件由引擎自动生成，通常在最终交付时删除，需要时用
\texttt{latexmk -c} 或手动清理。

\section{本手册的结构}

本手册按照"先总览、再入门、后原理"的顺序组织：

\begin{itemize}
  \item 第一章（本章）：介绍功能与模块划分；
  \item 第二章：从最小示例出发，给出完整的入门流程；
  \item 第三章：版面与排版的数学原理与公式推导；
  \item 第四章：定理环境的内部设计与算法实现；
  \item 第五章：编译管线与文档类的内部自动化机制；
  \item 第六章：进阶定制方法；
  \item 附录 A：文档类选项速查表。
\end{itemize}

阅读完前两章即可上手使用；若希望深入理解模板的工作原理，甚至为其
贡献代码，请重点阅读第三、四、五章。

\section{历史与版本演进}

ElegantBook 自 2013 年发布以来经历了多次大版本重构，
理解其演进脉络有助于读懂网上的旧教程：

\begin{table}[htbp]
  \centering
  \caption{ElegantBook 主要版本演进}\label{tab:history}
  \begin{tabular}{llp{7.6cm}}
    \toprule
    版本 & 时间 & 主要变化 \\
    \midrule
    v1.x & 2013--2018 & 面向理工科笔记的简单书籍模板，
          提供封面与基础定理环境 \\
    v2   & 2018      & 引入主题色（绿/青/蓝）与
          \texttt{fancy} 版式，封面可插图 \\
    v3   & 2020      & 用 tcolorbox 重构全部定理环境，
          支持跨页断行的彩色盒子 \\
    v4   & 2021 至今 & 选项键值化（kvoptions）、多语言
          （cn/en/fr/de/jp 等）、\texttt{simple} 简约模式、
          \texttt{device=pad} 平板版式、用户自定义定理环境
          \texttt{\textbackslash elegantnewtheorem} \\
    \bottomrule
  \end{tabular}
\end{table}

当前主线版本（本文对应 v4.6）在设计上遵循三条原则：
\textbf{选项即配置}（一切外观差异都由文档类选项表达）、
\textbf{用户零侵入}（封面、目录、页眉等自动化，
用户只需填元信息与正文）、\textbf{可扩展}（颜色、定理环境、
字体均可覆盖）。

\section{特性的代码示例}

下面把第一节的功能清单落到具体代码上，帮助读者建立直观印象。

\subsection{封面元信息}

\begin{lstlisting}
\title{数学分析讲义}
\subtitle{微积分与级数理论}
\author{张三 \and 李四}
\institute{某某大学数学学院}
\date{2026 年 8 月}
\version{1.0}
\extrainfo{仅供课堂使用}
\logo{figures/logo.pdf}     % 可选：封面左上角校徽
\cover{figures/cover.pdf}   % 可选：整页封面背景图
\end{lstlisting}

\verb|\author| 支持 \verb|\and| 分隔多个作者；
\verb|\cover| 若给出，将作为整页背景铺满封面。

\subsection{主题色}

\begin{lstlisting}
\documentclass[lang=cn, color=cyan]{elegantbook}   % 青色调
\documentclass[lang=cn, color=green]{elegantbook}  % 绿色调
\end{lstlisting}

也可以完全自定义（见第六章）。

\subsection{定理盒子}

\begin{lstlisting}
\begin{theorem}[欧拉公式]
对任意实数 $x$，有 $e^{ix}=\cos x+i\sin x$。
\end{theorem}

\begin{lemma}
设 $f$ 在闭区间 $[a,b]$ 上连续，则 $f$ 一致连续。
\end{lemma}
\end{lstlisting}

\subsection{参考文献}

\begin{lstlisting}
% 主文件导言区
\addbibresource{refs.bib}
% 正文中引用
\cite{knuth1984texbook}
% 文末输出
\printbibliography[heading=bibintoc, title={参考文献}]
\end{lstlisting}

\section{模块内部宏包明细}

第一节给出了宏观模块划分，下表列出每个模块实际加载的核心宏包
及其职责，便于读者判断"模板内部已经替我做了什么"：

\begin{table}[htbp]
  \centering
  \caption{模块与内部宏包明细}\label{tab:packages}
  \begin{tabular}{llp{6.5cm}}
    \toprule
    模块 & 宏包 & 职责 \\
    \midrule
    语言 & ctex & 中文标点、缩进、字体配置（\texttt{lang=cn} 时）\\
    \multirow{2}{*}{字体} & fontspec & XeLaTeX 下系统字体加载 \\
                          & newtxtext/newtxmath & Times 风格西文与数学字体 \\
    \multirow{2}{*}{版面} & geometry & 页边距（见第三章推导）\\
                          & titlesec, fancyhdr, tocloft & 标题、页眉页脚、目录样式 \\
    \multirow{2}{*}{数学} & amsmath 系列 & 公式环境、符号字体 \\
                          & bm, esint & 粗斜体、积分符号 \\
    \multirow{2}{*}{盒子} & tcolorbox & 定理盒子的绘制引擎 \\
                          & amsthm & 定理计数器与 proof 环境 \\
    图形图表 & tikz, graphicx, booktabs, multicol, multirow & 矢量图、插图、三线表 \\
    代码列表 & listings, fancyvrb, enumerate, enumitem & 代码高亮与列表 \\
    文献 & biblatex & 参考文献（后端 biber/bibtex）\\
    工具 & hyperref, xcolor, kvoptions, etoolbox & 链接、颜色、选项解析 \\
    \bottomrule
  \end{tabular}
\end{table}

\begin{remark}
表中大部分宏包读者无需自行加载；若需要\emph{覆盖}其配置，
应在主文件导言区于 \verb|\begin{document}| 之前重新设置
（例如用 \verb|\definecolor| 覆盖主题色）。
\end{remark}

\section{与其他排版方案的对比}

选型时常见的问题："我为什么不用 ctexbook 或直接堆宏包？"
下表给出粗粒度对比：

\begin{table}[htbp]
  \centering
  \caption{常见中文书籍排版方案对比}\label{tab:compare}
  \begin{tabular}{lcccc}
    \toprule
    方案 & 中文开箱 & 彩色定理盒子 & 自动封面 & 配置成本 \\
    \midrule
    \texttt{ctexbook} + 手写宏包 & 好 & 无，需自建 & 无，需自建 & 高 \\
    \texttt{tufte-book} & 一般 & 无 & 无 & 中 \\
    \texttt{memoir} & 一般 & 无 & 无 & 高 \\
    \textbf{ElegantBook} & 好 & 内置 & 内置 & 低 \\
    \bottomrule
  \end{tabular}
\end{table}

ElegantBook 的取舍是：牺牲一部分"从零控制的自由度"，换取
"两天出成品"的确定性。对需要深度定制的用户，模板同样开放了
颜色、字体、定理环境与页眉页脚的覆盖接口（见第六、七章
与附录~A）。

\section{运行环境要求}

\begin{itemize}
  \item \textbf{引擎}：XeLaTeX（中文必需）；纯英文文档亦可用
        pdfLaTeX；
  \item \textbf{发行版}：TeX~Live（推荐，2021 版以后）或 MiKTeX；
        安装 \texttt{elegantbook} 宏包：
\begin{lstlisting}
tlmgr install elegantbook
\end{lstlisting}
  \item \textbf{参考文献后端}：biber（默认）或 BibTeX；
  \item \textbf{字体}：中文排版需要系统中存在中文字体
        （Windows 自带宋体/黑体/楷体，macOS 自带苹方/宋体等）；
  \item \textbf{编辑工具}：TeXstudio、VS Code + LaTeX~Workshop、
        Overleaf（在线）均可；编辑器内配置"XeLaTeX 编译 +
        biber 参考文献"即可（见第二章）。
\end{itemize}

% ============================================================================
%  第二章  快速入门
% ============================================================================

\chapter{快速入门}

本章从一个最小示例出发，介绍从零开始编写一份 ElegantBook 中文文档的
完整流程：源文件怎么写、怎么编译、常用环境怎么用、出现问题怎么排错。
读者只需跟随本章操作即可产出第一份成品。

\section{一个最小示例}

下面是最小的完整源文件（假设文件名为 \texttt{main.tex}）：

\begin{lstlisting}
% main.tex —— 最小示例
\documentclass[lang=cn]{elegantbook}

\title{我的第一本书}
\author{张三}
\institute{某某大学}
\date{\today}

\begin{document}
\maketitle
\tableofcontents

\chapter{引言}
这是一本用 ElegantBook 排版的书籍。
\end{document}
\end{lstlisting}

将以上内容保存为 \texttt{main.tex} 后，依次执行第~\ref{sec:compile}~节给出的
四条命令，即可得到 \texttt{main.pdf}：其中封面、目录、章节标题的排版
全部由文档类自动完成。可以看到，用户只需要填写元信息与正文内容，
其余工作（封面绘制、页眉页脚、目录生成、超链接着色等）都由模板代劳。

\section{编译流程与引擎要求}\label{sec:compile}

ElegantBook 的中文模式需要 \textbf{XeLaTeX} 引擎（它原生支持 UTF-8 与系统
字体），因此请使用 XeLaTeX 编译；若使用 \texttt{lang=en} 的纯英文文档，
则 pdfLaTeX 亦可。推荐的完整编译命令为：

\begin{verbatim}
xelatex main.tex      % 第 1 遍：生成 .aux、.toc、.bcf 等辅助文件
biber   main          % 处理参考文献：读 .bcf，写 .bbl
xelatex main.tex      % 第 2 遍：引用参考文献并写入目录
xelatex main.tex      % 第 3 遍：收敛页码与交叉引用
\end{verbatim}

为什么需要多遍编译？因为 LaTeX 采用"两遍/多遍编译"模型：第一遍只知道
每个\emph{标签}（\verb|\label|）的位置，但页码与编号要等排版完成才能确定；
目录（\texttt{.toc}）、交叉引用（\texttt{.aux}）和参考文献（\texttt{.bbl}）
正是在各遍之间传递信息的文件。若正文或引用有变动，第三遍通常已足够收敛。

\begin{remark}
也可以使用 \texttt{latexmk} 自动完成上述流程：
\begin{verbatim}
latexmk -xelatex -interaction=nonstopmode main.tex
\end{verbatim}
该工具会监测源文件变化并自动重跑所需遍数，适合写作阶段使用。
\end{remark}

\section{文档类选项}

\verb|\documentclass| 的可选参数以键值对形式给出，最常用的组合如下：

\begin{lstlisting}
\documentclass[11pt,a4paper,lang=cn,color=blue,mode=fancy,
               device=normal,math=cm,titlestyle=hang]{elegantbook}
\end{lstlisting}

其中各选项的含义为：

\begin{description}
  \item[\texttt{lang=cn}] 语言环境，中文排版；可选 \texttt{en}（英文）、
        \texttt{fr}（法文）、\texttt{de}（德文）、\texttt{jp}（日文）等；
  \item[\texttt{color=blue}] 主题色，可选 \texttt{green}、\texttt{cyan}、
        \texttt{blue}、\texttt{gray}、\texttt{black}、\texttt{nocolor}；
  \item[\texttt{mode=fancy}] 版式，\texttt{fancy} 为华丽版（默认），
        \texttt{simple} 为简约版；
  \item[\texttt{device=normal}] 输出设备，\texttt{normal}（默认）或
        \texttt{pad}（平板优化）；
  \item[\texttt{math=cm}] 数学字体，\texttt{cm}（默认，Computer Modern）、
        \texttt{newtx}（Times 风格）或 \texttt{mtpro2}；
  \item[\texttt{titlestyle=hang}] 章节标题样式，\texttt{hang}（悬挂式，
        默认）或 \texttt{display}（通栏式）。
\end{description}

其余选项（如 \texttt{chinesefont}、\texttt{toc}、\texttt{thmcnt}、
\texttt{citestyle}、\texttt{bibend} 等）在附录~A 中给出完整速查表。

\section{封面与前置部分}

在正文之前，需要填写封面元信息。ElegantBook 提供的封面命令包括：

\begin{lstlisting}
\title{书名}                 % 必须
\subtitle{副标题}            % 可选
\author{作者}                % 必须
\institute{单位/学校}        % 可选
\date{2026 年 8 月}          % 可选，默认当天
\version{1.0}                % 可选，显示在封面
\extrainfo{副信息}           % 可选
\logo{figures/logo.pdf}      % 可选，封面校徽
\cover{figures/cover.pdf}    % 可选，整页封面图
\keywords{关键词}            % 可选
\end{lstlisting}

然后在 \verb|\begin{document}| 之后依次调用
\verb|\maketitle|、\verb|\tableofcontents|，并用
\verb|\frontmatter|/\verb|\mainmatter| 划分前置部分与正文：

\begin{lstlisting}
\begin{document}
\maketitle
\frontmatter
\tableofcontents
\mainmatter
\include{chapters/ch01_overview}
\include{chapters/ch02_quickstart}
...
\end{document}
\end{lstlisting}

其中 \verb|\include{chapters/xxx}| 用于按章分文件组织正文，
这是大型文档推荐的写作方式：每一章一个 \texttt{.tex} 文件，
正文使用 \verb|\chapter{}| 开头，编译时由主文件统一导入。

\section{正文常用环境}

\subsection{定理与证明}

ElegantBook 提供了完整的定理盒子体系，用法与 \texttt{amsthm} 一致，
只是外观被包装成了带主题色的圆角盒子：

\begin{lstlisting}
\begin{definition}[收敛]
设 $\{x_n\}$ 是实数序列，若存在 $a\in\mathbb{R}$，使得对任意
$\varepsilon>0$，存在 $N\in\mathbb{N}$，当 $n>N$ 时有
$|x_n-a|<\varepsilon$，则称 $\{x_n\}$ 收敛于 $a$。
\end{definition}

\begin{theorem}[夹逼定理]
若 $x_n\leqslant y_n\leqslant z_n$ 且 $\lim x_n=\lim z_n=a$，
则 $\lim y_n=a$。
\end{theorem}

\begin{proof}
由定义直接验证即可。
\end{proof}
\end{lstlisting}

实际效果如本节下面的例~\ref{ex:converge} 所示（定理环境的具体实现原理
见第四章）：

\begin{definition}[收敛]\label{ex:converge}
设 $\{x_n\}$ 是实数序列，若存在 $a\in\mathbb{R}$，使得对任意
$\varepsilon>0$，存在 $N\in\mathbb{N}$，当 $n>N$ 时有
$|x_n-a|<\varepsilon$，则称 $\{x_n\}$ 收敛于 $a$，记作
$\lim\limits_{n\to\infty}x_n=a$。
\end{definition}

\subsection{数学公式}

文档类已加载 \texttt{amsmath} 系列宏包，行内公式与多行公式均可直接使用：

\begin{lstlisting}
\begin{equation}
  e^{i\pi}+1=0 \label{eq:euler}
\end{equation}

\begin{align}
  (a+b)^2 &= a^2+2ab+b^2,\\
  (a-b)^2 &= a^2-2ab+b^2.
\end{align}
\end{lstlisting}

例如欧拉公式（\eqref{eq:euler}）被许多人称为数学中最美的公式之一：
\begin{equation}
  e^{i\pi}+1=0 \label{eq:euler}
\end{equation}

\subsection{插图}

插图既可以使用外部图片（\texttt{graphicx}），也可以直接用 \texttt{tikz}
绘制矢量图，后者无需额外文件、更利于保持文档自包含：

\begin{lstlisting}
\begin{figure}[htbp]
  \centering
  \begin{tikzpicture}
    \draw[fill=structurecolor!20] (0,0) circle (1);
    \draw[->,thick] (0,0) -- (0.707,0.707);
    \node at (0.4,0.9) {\scriptsize $\theta$};
  \end{tikzpicture}
  \caption{一个单位圆}\label{fig:unitcircle}
\end{figure}
\end{lstlisting}

引用插图时可以使用文档类提供的 \verb|\figref|（或
\verb|\tabref| 引用表格），例如"见图~\figref{fig:unitcircle}"：

\begin{figure}[htbp]
  \centering
  \begin{tikzpicture}
    \draw[fill=structurecolor!20] (0,0) circle (1);
    \draw[->,thick] (0,0) -- (0.707,0.707);
    \node at (0.42,0.55) {\scriptsize $\theta$};
    \draw[<->] (1.2,0) arc (0:45:1.2);
  \end{tikzpicture}
  \caption{一个单位圆}\label{fig:unitcircle}
\end{figure}

\subsection{表格}

三线表使用 \texttt{booktabs} 的 \verb|\toprule|、
\verb|\midrule|、\verb|\bottomrule|：

\begin{lstlisting}
\begin{table}[htbp]
  \centering
  \caption{示例数据}\label{tab:data}
  \begin{tabular}{lcc}
    \toprule
    变量 & 均值 & 方差 \\
    \midrule
    $x$ & 0.0 & 1.0 \\
    $y$ & 1.5 & 0.4 \\
    \bottomrule
  \end{tabular}
\end{table}
\end{lstlisting}

\subsection{代码列表}

源代码环境由内置的 \texttt{listings} 配置提供，
本手册中的全部示例代码均使用 \texttt{lstlisting} 环境排版。

\subsection{列表}

无序、有序、描述三种列表均可用，且项目符号被替换为主题色圆点：

\begin{lstlisting}
\begin{itemize}
  \item 无序列表项
\end{itemize}
\begin{enumerate}
  \item 有序列表项
\end{enumerate}
\begin{description}
  \item[术语] 解释文字。
\end{description}
\end{lstlisting}

\section{参考文献的使用}

参考文献数据库写在 \texttt{refs.bib} 中，主文件前言用
\verb|\addbibresource{refs.bib}| 引入：

\begin{lstlisting}
% refs.bib
@book{lamport1994latex,
  author    = {Lamport, Leslie},
  title     = {LaTeX: A Document Preparation System},
  edition   = {2nd},
  publisher = {Addison-Wesley},
  year      = {1994},
}
\end{lstlisting}

正文中用 \verb|\cite{...}| 引用，文末用
\verb|\printbibliography| 输出文献列表。例如
\cite{lamport1994latex} 是 LaTeX 的经典入门著作。注意：引入了
文献后，编译时必须在两次 \texttt{xelatex} 之间运行 \texttt{biber}，
否则引用位置会出现问号。

\section{常见问题与排错}

\begin{description}
  \item[\texttt{elegantbook.cls} 找不到] 说明发行版未安装该文档类。
        解决方法：运行 \texttt{tlmgr install elegantbook}（TeX~Live），
        或者把 \texttt{elegantbook.cls} 直接复制到源文件所在目录（推荐，
        本手册即采用此方式保证自包含）；
  \item[\texttt{File xxx.sty not found}] 缺少宏包，用
        \texttt{tlmgr install xxx} 安装，或检查是否误用了 pdfLaTeX
        编译中文文档；
  \item[编译后中文乱码/缺字] 确认源文件为 UTF-8 编码，并确认系统中
        存在中文字体（Windows 自带宋体/黑体/楷体）；
  \item[引用显示为问号] 交叉引用需多编一次；参考文献需运行
        \texttt{biber}；
  \item[\texttt{biber} 报错] 检查主文件选项 \texttt{bibend} 与后端
        是否一致（默认 \texttt{biber}，若用 BibTeX 后端则
        需 \texttt{bibend=bibtex}）。
\end{description}

\section{最小示例逐行讲解}

把第一节的最小示例逐行拆开，说明每一行的作用：

\begin{table}[htbp]
  \centering
  \caption{最小示例逐行注释}\label{tab:linebyline}
  \begin{tabular}{lp{9.5cm}}
    \toprule
    代码 & 作用 \\
    \midrule
    \texttt{\textbackslash documentclass[lang=cn]\{elegantbook\}}
      & 声明使用 ElegantBook 文档类，启用中文模式；
        之后才可编写正文 \\
    \texttt{\textbackslash title\{...\}} 等四行
      & 填写封面元信息：书名、作者、单位、日期 \\
    \texttt{\textbackslash begin\{document\}}
      & 正文开始（此前为导言区，只能放设置命令）\\
    \texttt{\textbackslash maketitle}
      & 按元信息自动排版整页封面并输出 \\
    \texttt{\textbackslash tableofcontents}
      & 输出目录（内容来自 \texttt{.toc} 文件）\\
    \texttt{\textbackslash chapter\{引言\}}
      & 开启一章；章标题自动进入目录与页眉 \\
    \texttt{\textbackslash end\{document\}}
      & 正文结束；\texttt{\textbackslash printbibliography} 等
        收尾命令应在它之前 \\
    \bottomrule
  \end{tabular}
\end{table}

注意导言区（\verb|\begin{document}| 之前）只能放\textbf{设置类}命令，
正文内容必须放在文档环境内——这是 LaTeX 最基本的语法边界。

\section{辅助文件详解}

编译过程中产生的辅助文件各有分工，了解它们有助于排查问题：

\begin{table}[htbp]
  \centering
  \caption{辅助文件一览}\label{tab:auxfiles}
  \begin{tabular}{llp{7.5cm}}
    \toprule
    文件 & 生产者 & 内容与作用 \\
    \midrule
    \texttt{.aux} & xelatex & 交叉引用表：\texttt{\textbackslash label}
          与 \texttt{\textbackslash ref} 的信息交换 \\
    \texttt{.toc} & xelatex & 目录条目，由 \texttt{\textbackslash
          chapter} 等写入 \\
    \texttt{.out} & hyperref & PDF 书签树 \\
    \texttt{.bcf} & biblatex & 文献引用记录（第一次 xelatex 写入）\\
    \texttt{.bbl} & biber & 文献列表排版结果，xelatex 直接读入 \\
    \texttt{.blg} & biber & biber 运行日志（报错时先看它）\\
    \texttt{.log} & xelatex & 编译日志（报错时的主要线索）\\
    \texttt{.synctex.gz} & xelatex & 源码与 PDF 双向定位（编辑器用）\\
    \bottomrule
  \end{tabular}
\end{table}

这些文件都可放心删除——下次编译会自动重建。唯一需要保留的是
\texttt{main.tex}、\texttt{chapters/} 与 \texttt{refs.bib}。

\section{三种编译方式}

\subsection{手动命令（最透明）}

适合理解原理与排查问题：
\begin{verbatim}
xelatex main.tex && biber main && xelatex main.tex && xelatex main.tex
\end{verbatim}

\subsection{latexmk（写作期推荐）}

latexmk 自动判断需要重跑几遍。在项目根目录放一个
\texttt{.latexmkrc}（Windows 上亦可），内容为：
\begin{lstlisting}
$pdf_mode     = 5;      # 用 xelatex 编译
$bibtex_use   = 2;      # 需要时自动运行参考文献工具
$biber        = 'biber %O %S';   # 指定 biber 后端
\end{lstlisting}
之后只需一条命令：
\begin{verbatim}
latexmk main.tex
\end{verbatim}

\subsection{编辑器集成}

\begin{description}
  \item[TeXstudio] 选项 $\rightarrow$ 构建 $\rightarrow$ 默认编译器
        选 \texttt{XeLaTeX}；文献工具选 \texttt{biber}；
  \item[VS Code + LaTeX Workshop] 在 \texttt{settings.json} 中把
        \texttt{latex-workshop.latex.recipes} 第一条配成
        \texttt{xelatex $\rightarrow$ biber $\rightarrow$ xelatex $\rightarrow$
        xelatex}；
  \item[Overleaf] 菜单 $\rightarrow$ 编译器选 \texttt{XeLaTeX}；
        Overleaf 会自动调用 biber。
\end{description}

\section{常用环境速查}

\subsection{数学公式环境}

文档类已加载 amsmath，常用环境如下：

\begin{table}[htbp]
  \centering
  \caption{常用数学环境}\label{tab:mathenv}
  \begin{tabular}{lll}
    \toprule
    环境 & 用途 & 编号 \\
    \midrule
    \texttt{equation} / \texttt{equation*} & 单行公式 & 有/无 \\
    \texttt{align} / \texttt{align*} & 多行按 \texttt{\&} 对齐 & 有/无 \\
    \texttt{gather} / \texttt{gather*} & 多行居中 & 有/无 \\
    \texttt{multline} & 超长公式换行（首行左、末行右）& 有 \\
    \texttt{cases} & 分段函数 & 嵌套使用 \\
    \texttt{split} & 多行共享一个编号 & 共享 \\
    \texttt{pmatrix}/\texttt{bmatrix} & 圆/方括号矩阵 & --- \\
    \bottomrule
  \end{tabular}
\end{table}

例如分段函数与矩阵：

\begin{lstlisting}
\[
  f(x)=\begin{cases}
    x^2, & x\geqslant 0,\\
    -x,  & x<0,
  \end{cases}
  \qquad
  A=\begin{pmatrix} a & b \\ c & d \end{pmatrix}
\]
\end{lstlisting}

渲染结果为
\[
  f(x)=\begin{cases}
    x^2, & x\geqslant 0,\\
    -x,  & x<0,
  \end{cases}
  \qquad
  A=\begin{pmatrix} a & b \\ c & d \end{pmatrix}.
\]

\subsection{其他常用环境}

\begin{description}
  \item[\texttt{itemize}/\texttt{enumerate}/\texttt{description}]
        无序/有序/描述列表，项目符号为主题色圆点；
  \item[\texttt{lstlisting}] 代码环境，支持语法高亮；
  \item[\texttt{verbatim}] 原样输出（不处理命令）；
  \item[\texttt{quote}/\texttt{quotation}] 短/长引文；
  \item[\texttt{abstract}] 摘要（文档类自带样式）；
  \item[\texttt{introduction}] 章前引言框（双栏要点列表）。
\end{description}

\section{标签与交叉引用规范}

\verb|\label| 与 \verb|\ref| 是文档内引用的基础。推荐的
标签命名规范是"类型前缀 + 冒号 + 语义名"：

\begin{table}[htbp]
  \centering
  \caption{标签前缀约定}\label{tab:labels}
  \begin{tabular}{lll}
    \toprule
    前缀 & 对象 & 引用命令 \\
    \midrule
    \texttt{fig:} & 图 & \verb|\ref|、\verb|\figref| \\
    \texttt{tab:} & 表 & \verb|\ref|、\verb|\tabref| \\
    \texttt{eq:} & 公式 & \verb|\eqref|（自动加括号）\\
    \texttt{sec:}/\texttt{ch:} & 节/章 & \verb|\ref|、\verb|\autoref| \\
    \texttt{thm:}/\texttt{def:}/\texttt{lem:} & 定理/定义/引理 & \verb|\ref| \\
    \texttt{alg:} & 算法 & \verb|\ref| \\
    \bottomrule
  \end{tabular}
\end{table}

\verb|\label| 必须放在产生编号的环境\emph{内部}（例如
\verb|\begin{equation}...\label{eq:x}\end{equation}|），
它记录的是"当前计数器的值"，放在外面会捕获错误编号。

\section{章节层级与编号深度}

ElegantBook 继承 \texttt{book} 类的层级结构：
\begin{equation*}
  \texttt{part} \succ \texttt{chapter} \succ \texttt{section}
  \succ \texttt{subsection} \succ \texttt{subsubsection}
  \succ \texttt{paragraph} \succ \texttt{subparagraph}.
\end{equation*}
默认编号深度 \texttt{secnumdepth=2}（即编号到
\texttt{subsection}），目录深度 \texttt{tocdepth=2}。
需要更深时可调整：
\begin{lstlisting}
\setcounter{secnumdepth}{3}   % subsubsection 也编号
\setcounter{tocdepth}{3}      % 目录包含 subsubsection
\end{lstlisting}

\section{一个完整的示例章}

综合以上内容，下面是一章完整的写法模板
（保存为 \texttt{chapters/ch01\_demo.tex} 并被主文件
\verb|\include|）：

\begin{lstlisting}
% chapters/ch01_demo.tex
\chapter{示例章}

\section{引入}
本节对应 \ref{sec:model} 的模型。

\begin{definition}[导数]\label{def:deriv}
设 $f$ 在 $x_0$ 附近有定义，若极限
\[ f'(x_0)=\lim_{h\to0}\frac{f(x_0+h)-f(x_0)}{h} \]
存在，则称 $f$ 在 $x_0$ 处可导。
\end{definition}

\begin{theorem}[可导必连续]
若 $f$ 在 $x_0$ 处可导，则 $f$ 在 $x_0$ 处连续。
\end{theorem}
\begin{proof}
由 $f'(x_0)$ 存在直接推出
$\lim_{x\to x_0}\big(f(x)-f(x_0)\big)=0$。
\end{proof}

\section{模型}\label{sec:model}
线性近似（\eqref{eq:linear}）是微积分的基本思想：
\begin{equation}
  f(x)\approx f(x_0)+f'(x_0)(x-x_0). \label{eq:linear}
\end{equation}

\begin{figure}[htbp]
  \centering
  \begin{tikzpicture}
    \draw[->] (-0.5,0) -- (3,0) node[right] {$x$};
    \draw[->] (0,-0.5) -- (0,3) node[above] {$y$};
    \draw[thick, structurecolor] plot[domain=0:2.5]
      (\x, {0.5*\x*\x});
    \draw[thick, second] (0.5,0.125) -- (2.5,2.125);
  \end{tikzpicture}
  \caption{割线逼近切线}\label{fig:approx}
\end{figure}

图像见图~\figref{fig:approx}，理论见 \cite{knuth1984texbook}。
\end{lstlisting}

按此模板组织，一本书的每个章节彼此独立，便于多人协作与单独编译调试。

...
\end{document}
\end{lstlisting}

其中 \verb|\include{chapters/xxx}| 用于按章分文件组织正文，
这是大型文档推荐的写作方式：每一章一个 \texttt{.tex} 文件，
正文使用 \verb|\chapter{}| 开头，编译时由主文件统一导入。

\section{正文常用环境}

\subsection{定理与证明}

ElegantBook 提供了完整的定理盒子体系，用法与 \texttt{amsthm} 一致，
只是外观被包装成了带主题色的圆角盒子：

\begin{lstlisting}
\begin{definition}[收敛]
设 $\{x_n\}$ 是实数序列，若存在 $a\in\mathbb{R}$，使得对任意
$\varepsilon>0$，存在 $N\in\mathbb{N}$，当 $n>N$ 时有
$|x_n-a|<\varepsilon$，则称 $\{x_n\}$ 收敛于 $a$。
\end{definition}

\begin{theorem}[夹逼定理]
若 $x_n\leqslant y_n\leqslant z_n$ 且 $\lim x_n=\lim z_n=a$，
则 $\lim y_n=a$。
\end{theorem}

\begin{proof}
由定义直接验证即可。
\end{proof}
\end{lstlisting}

实际效果如本节下面的例~\ref{ex:converge} 所示（定理环境的具体实现原理
见第四章）：

\begin{definition}[收敛]\label{ex:converge}
设 $\{x_n\}$ 是实数序列，若存在 $a\in\mathbb{R}$，使得对任意
$\varepsilon>0$，存在 $N\in\mathbb{N}$，当 $n>N$ 时有
$|x_n-a|<\varepsilon$，则称 $\{x_n\}$ 收敛于 $a$，记作
$\lim\limits_{n\to\infty}x_n=a$。
\end{definition}

\subsection{数学公式}

文档类已加载 \texttt{amsmath} 系列宏包，行内公式与多行公式均可直接使用：

\begin{lstlisting}
\begin{equation}
  e^{i\pi}+1=0 \label{eq:euler}
\end{equation}

\begin{align}
  (a+b)^2 &= a^2+2ab+b^2,\\
  (a-b)^2 &= a^2-2ab+b^2.
\end{align}
\end{lstlisting}

例如欧拉公式（\eqref{eq:euler}）被许多人称为数学中最美的公式之一：
\begin{equation}
  e^{i\pi}+1=0 \label{eq:euler}
\end{equation}

\subsection{插图}

插图既可以使用外部图片（\texttt{graphicx}），也可以直接用 \texttt{tikz}
绘制矢量图，后者无需额外文件、更利于保持文档自包含：

\begin{lstlisting}
\begin{figure}[htbp]
  \centering
  \begin{tikzpicture}
    \draw[fill=structurecolor!20] (0,0) circle (1);
    \draw[->,thick] (0,0) -- (0.707,0.707);
    \node at (0.4,0.9) {\scriptsize $\theta$};
  \end{tikzpicture}
  \caption{一个单位圆}\label{fig:unitcircle}
\end{figure}
\end{lstlisting}

引用插图时可以使用文档类提供的 \verb|\figref|（或
\verb|\tabref| 引用表格），例如"见图~\figref{fig:unitcircle}"：

\begin{figure}[htbp]
  \centering
  \begin{tikzpicture}
    \draw[fill=structurecolor!20] (0,0) circle (1);
    \draw[->,thick] (0,0) -- (0.707,0.707);
    \node at (0.42,0.55) {\scriptsize $\theta$};
    \draw[<->] (1.2,0) arc (0:45:1.2);
  \end{tikzpicture}
  \caption{一个单位圆}\label{fig:unitcircle}
\end{figure}

\subsection{表格}

三线表使用 \texttt{booktabs} 的 \verb|\toprule|、
\verb|\midrule|、\verb|\bottomrule|：

\begin{lstlisting}
\begin{table}[htbp]
  \centering
  \caption{示例数据}\label{tab:data}
  \begin{tabular}{lcc}
    \toprule
    变量 & 均值 & 方差 \\
    \midrule
    $x$ & 0.0 & 1.0 \\
    $y$ & 1.5 & 0.4 \\
    \bottomrule
  \end{tabular}
\end{table}
\end{lstlisting}

\subsection{代码列表}

源代码环境由内置的 \texttt{listings} 配置提供，
本手册中的全部示例代码均使用 \texttt{lstlisting} 环境排版。

\subsection{列表}

无序、有序、描述三种列表均可用，且项目符号被替换为主题色圆点：

\begin{lstlisting}
\begin{itemize}
  \item 无序列表项
\end{itemize}
\begin{enumerate}
  \item 有序列表项
\end{enumerate}
\begin{description}
  \item[术语] 解释文字。
\end{description}
\end{lstlisting}

\section{参考文献的使用}

参考文献数据库写在 \texttt{refs.bib} 中，主文件前言用
\verb|\addbibresource{refs.bib}| 引入：

\begin{lstlisting}
% refs.bib
@book{lamport1994latex,
  author    = {Lamport, Leslie},
  title     = {LaTeX: A Document Preparation System},
  edition   = {2nd},
  publisher = {Addison-Wesley},
  year      = {1994},
}
\end{lstlisting}

正文中用 \verb|\cite{...}| 引用，文末用
\verb|\printbibliography| 输出文献列表。例如
\cite{lamport1994latex} 是 LaTeX 的经典入门著作。注意：引入了
文献后，编译时必须在两次 \texttt{xelatex} 之间运行 \texttt{biber}，
否则引用位置会出现问号。

\section{常见问题与排错}

\begin{description}
  \item[\texttt{elegantbook.cls} 找不到] 说明发行版未安装该文档类。
        解决方法：运行 \texttt{tlmgr install elegantbook}（TeX~Live），
        或者把 \texttt{elegantbook.cls} 直接复制到源文件所在目录（推荐，
        本手册即采用此方式保证自包含）；
  \item[\texttt{File xxx.sty not found}] 缺少宏包，用
        \texttt{tlmgr install xxx} 安装，或检查是否误用了 pdfLaTeX
        编译中文文档；
  \item[编译后中文乱码/缺字] 确认源文件为 UTF-8 编码，并确认系统中
        存在中文字体（Windows 自带宋体/黑体/楷体）；
  \item[引用显示为问号] 交叉引用需多编一次；参考文献需运行
        \texttt{biber}；
  \item[\texttt{biber} 报错] 检查主文件选项 \texttt{bibend} 与后端
        是否一致（默认 \texttt{biber}，若用 BibTeX 后端则
        需 \texttt{bibend=bibtex}）。
\end{description}

\section{最小示例逐行讲解}

把第一节的最小示例逐行拆开，说明每一行的作用：

\begin{table}[htbp]
  \centering
  \caption{最小示例逐行注释}\label{tab:linebyline}
  \begin{tabular}{lp{9.5cm}}
    \toprule
    代码 & 作用 \\
    \midrule
    \texttt{\textbackslash documentclass[lang=cn]\{elegantbook\}}
      & 声明使用 ElegantBook 文档类，启用中文模式；
        之后才可编写正文 \\
    \texttt{\textbackslash title\{...\}} 等四行
      & 填写封面元信息：书名、作者、单位、日期 \\
    \texttt{\textbackslash begin\{document\}}
      & 正文开始（此前为导言区，只能放设置命令）\\
    \texttt{\textbackslash maketitle}
      & 按元信息自动排版整页封面并输出 \\
    \texttt{\textbackslash tableofcontents}
      & 输出目录（内容来自 \texttt{.toc} 文件）\\
    \texttt{\textbackslash chapter\{引言\}}
      & 开启一章；章标题自动进入目录与页眉 \\
    \texttt{\textbackslash end\{document\}}
      & 正文结束；\texttt{\textbackslash printbibliography} 等
        收尾命令应在它之前 \\
    \bottomrule
  \end{tabular}
\end{table}

注意导言区（\verb|\begin{document}| 之前）只能放\textbf{设置类}命令，
正文内容必须放在文档环境内——这是 LaTeX 最基本的语法边界。

\section{辅助文件详解}

编译过程中产生的辅助文件各有分工，了解它们有助于排查问题：

\begin{table}[htbp]
  \centering
  \caption{辅助文件一览}\label{tab:auxfiles}
  \begin{tabular}{llp{7.5cm}}
    \toprule
    文件 & 生产者 & 内容与作用 \\
    \midrule
    \texttt{.aux} & xelatex & 交叉引用表：\texttt{\textbackslash label}
          与 \texttt{\textbackslash ref} 的信息交换 \\
    \texttt{.toc} & xelatex & 目录条目，由 \texttt{\textbackslash
          chapter} 等写入 \\
    \texttt{.out} & hyperref & PDF 书签树 \\
    \texttt{.bcf} & biblatex & 文献引用记录（第一次 xelatex 写入）\\
    \texttt{.bbl} & biber & 文献列表排版结果，xelatex 直接读入 \\
    \texttt{.blg} & biber & biber 运行日志（报错时先看它）\\
    \texttt{.log} & xelatex & 编译日志（报错时的主要线索）\\
    \texttt{.synctex.gz} & xelatex & 源码与 PDF 双向定位（编辑器用）\\
    \bottomrule
  \end{tabular}
\end{table}

这些文件都可放心删除——下次编译会自动重建。唯一需要保留的是
\texttt{main.tex}、\texttt{chapters/} 与 \texttt{refs.bib}。

\section{三种编译方式}

\subsection{手动命令（最透明）}

适合理解原理与排查问题：
\begin{verbatim}
xelatex main.tex && biber main && xelatex main.tex && xelatex main.tex
\end{verbatim}

\subsection{latexmk（写作期推荐）}

latexmk 自动判断需要重跑几遍。在项目根目录放一个
\texttt{.latexmkrc}（Windows 上亦可），内容为：
\begin{lstlisting}
$pdf_mode     = 5;      # 用 xelatex 编译
$bibtex_use   = 2;      # 需要时自动运行参考文献工具
$biber        = 'biber %O %S';   # 指定 biber 后端
\end{lstlisting}
之后只需一条命令：
\begin{verbatim}
latexmk main.tex
\end{verbatim}

\subsection{编辑器集成}

\begin{description}
  \item[TeXstudio] 选项 $\rightarrow$ 构建 $\rightarrow$ 默认编译器
        选 \texttt{XeLaTeX}；文献工具选 \texttt{biber}；
  \item[VS Code + LaTeX Workshop] 在 \texttt{settings.json} 中把
        \texttt{latex-workshop.latex.recipes} 第一条配成
        \texttt{xelatex $\rightarrow$ biber $\rightarrow$ xelatex $\rightarrow$
        xelatex}；
  \item[Overleaf] 菜单 $\rightarrow$ 编译器选 \texttt{XeLaTeX}；
        Overleaf 会自动调用 biber。
\end{description}

\section{常用环境速查}

\subsection{数学公式环境}

文档类已加载 amsmath，常用环境如下：

\begin{table}[htbp]
  \centering
  \caption{常用数学环境}\label{tab:mathenv}
  \begin{tabular}{lll}
    \toprule
    环境 & 用途 & 编号 \\
    \midrule
    \texttt{equation} / \texttt{equation*} & 单行公式 & 有/无 \\
    \texttt{align} / \texttt{align*} & 多行按 \texttt{\&} 对齐 & 有/无 \\
    \texttt{gather} / \texttt{gather*} & 多行居中 & 有/无 \\
    \texttt{multline} & 超长公式换行（首行左、末行右）& 有 \\
    \texttt{cases} & 分段函数 & 嵌套使用 \\
    \texttt{split} & 多行共享一个编号 & 共享 \\
    \texttt{pmatrix}/\texttt{bmatrix} & 圆/方括号矩阵 & --- \\
    \bottomrule
  \end{tabular}
\end{table}

例如分段函数与矩阵：

\begin{lstlisting}
\[
  f(x)=\begin{cases}
    x^2, & x\geqslant 0,\\
    -x,  & x<0,
  \end{cases}
  \qquad
  A=\begin{pmatrix} a & b \\ c & d \end{pmatrix}
\]
\end{lstlisting}

渲染结果为
\[
  f(x)=\begin{cases}
    x^2, & x\geqslant 0,\\
    -x,  & x<0,
  \end{cases}
  \qquad
  A=\begin{pmatrix} a & b \\ c & d \end{pmatrix}.
\]

\subsection{其他常用环境}

\begin{description}
  \item[\texttt{itemize}/\texttt{enumerate}/\texttt{description}]
        无序/有序/描述列表，项目符号为主题色圆点；
  \item[\texttt{lstlisting}] 代码环境，支持语法高亮；
  \item[\texttt{verbatim}] 原样输出（不处理命令）；
  \item[\texttt{quote}/\texttt{quotation}] 短/长引文；
  \item[\texttt{abstract}] 摘要（文档类自带样式）；
  \item[\texttt{introduction}] 章前引言框（双栏要点列表）。
\end{description}

\section{标签与交叉引用规范}

\verb|\label| 与 \verb|\ref| 是文档内引用的基础。推荐的
标签命名规范是"类型前缀 + 冒号 + 语义名"：

\begin{table}[htbp]
  \centering
  \caption{标签前缀约定}\label{tab:labels}
  \begin{tabular}{lll}
    \toprule
    前缀 & 对象 & 引用命令 \\
    \midrule
    \texttt{fig:} & 图 & \verb|\ref|、\verb|\figref| \\
    \texttt{tab:} & 表 & \verb|\ref|、\verb|\tabref| \\
    \texttt{eq:} & 公式 & \verb|\eqref|（自动加括号）\\
    \texttt{sec:}/\texttt{ch:} & 节/章 & \verb|\ref|、\verb|\autoref| \\
    \texttt{thm:}/\texttt{def:}/\texttt{lem:} & 定理/定义/引理 & \verb|\ref| \\
    \texttt{alg:} & 算法 & \verb|\ref| \\
    \bottomrule
  \end{tabular}
\end{table}

\verb|\label| 必须放在产生编号的环境\emph{内部}（例如
\verb|\begin{equation}...\label{eq:x}\end{equation}|），
它记录的是"当前计数器的值"，放在外面会捕获错误编号。

\section{章节层级与编号深度}

ElegantBook 继承 \texttt{book} 类的层级结构：
\begin{equation*}
  \texttt{part} \succ \texttt{chapter} \succ \texttt{section}
  \succ \texttt{subsection} \succ \texttt{subsubsection}
  \succ \texttt{paragraph} \succ \texttt{subparagraph}.
\end{equation*}
默认编号深度 \texttt{secnumdepth=2}（即编号到
\texttt{subsection}），目录深度 \texttt{tocdepth=2}。
需要更深时可调整：
\begin{lstlisting}
\setcounter{secnumdepth}{3}   % subsubsection 也编号
\setcounter{tocdepth}{3}      % 目录包含 subsubsection
\end{lstlisting}

\section{一个完整的示例章}

综合以上内容，下面是一章完整的写法模板
（保存为 \texttt{chapters/ch01\_demo.tex} 并被主文件
\verb|\include|）：

\begin{lstlisting}
% chapters/ch01_demo.tex
\chapter{示例章}

\section{引入}
本节对应 \ref{sec:model} 的模型。

\begin{definition}[导数]\label{def:deriv}
设 $f$ 在 $x_0$ 附近有定义，若极限
\[ f'(x_0)=\lim_{h\to0}\frac{f(x_0+h)-f(x_0)}{h} \]
存在，则称 $f$ 在 $x_0$ 处可导。
\end{definition}

\begin{theorem}[可导必连续]
若 $f$ 在 $x_0$ 处可导，则 $f$ 在 $x_0$ 处连续。
\end{theorem}
\begin{proof}
由 $f'(x_0)$ 存在直接推出
$\lim_{x\to x_0}\big(f(x)-f(x_0)\big)=0$。
\end{proof}

\section{模型}\label{sec:model}
线性近似（\eqref{eq:linear}）是微积分的基本思想：
\begin{equation}
  f(x)\approx f(x_0)+f'(x_0)(x-x_0). \label{eq:linear}
\end{equation}

\begin{figure}[htbp]
  \centering
  \begin{tikzpicture}
    \draw[->] (-0.5,0) -- (3,0) node[right] {$x$};
    \draw[->] (0,-0.5) -- (0,3) node[above] {$y$};
    \draw[thick, structurecolor] plot[domain=0:2.5]
      (\x, {0.5*\x*\x});
    \draw[thick, second] (0.5,0.125) -- (2.5,2.125);
  \end{tikzpicture}
  \caption{割线逼近切线}\label{fig:approx}
\end{figure}

图像见图~\figref{fig:approx}，理论见 \cite{knuth1984texbook}。
\end{lstlisting}

按此模板组织，一本书的每个章节彼此独立，便于多人协作与单独编译调试。

...
\end{document}
\end{lstlisting}

其中 \verb|\include{chapters/xxx}| 用于按章分文件组织正文，
这是大型文档推荐的写作方式：每一章一个 \texttt{.tex} 文件，
正文使用 \verb|\chapter{}| 开头，编译时由主文件统一导入。

\section{正文常用环境}

\subsection{定理与证明}

ElegantBook 提供了完整的定理盒子体系，用法与 \texttt{amsthm} 一致，
只是外观被包装成了带主题色的圆角盒子：

\begin{lstlisting}
\begin{definition}[收敛]
设 $\{x_n\}$ 是实数序列，若存在 $a\in\mathbb{R}$，使得对任意
$\varepsilon>0$，存在 $N\in\mathbb{N}$，当 $n>N$ 时有
$|x_n-a|<\varepsilon$，则称 $\{x_n\}$ 收敛于 $a$。
\end{definition}

\begin{theorem}[夹逼定理]
若 $x_n\leqslant y_n\leqslant z_n$ 且 $\lim x_n=\lim z_n=a$，
则 $\lim y_n=a$。
\end{theorem}

\begin{proof}
由定义直接验证即可。
\end{proof}
\end{lstlisting}

实际效果如本节下面的例~\ref{ex:converge} 所示（定理环境的具体实现原理
见第四章）：

\begin{definition}[收敛]\label{ex:converge}
设 $\{x_n\}$ 是实数序列，若存在 $a\in\mathbb{R}$，使得对任意
$\varepsilon>0$，存在 $N\in\mathbb{N}$，当 $n>N$ 时有
$|x_n-a|<\varepsilon$，则称 $\{x_n\}$ 收敛于 $a$，记作
$\lim\limits_{n\to\infty}x_n=a$。
\end{definition}

\subsection{数学公式}

文档类已加载 \texttt{amsmath} 系列宏包，行内公式与多行公式均可直接使用：

\begin{lstlisting}
\begin{equation}
  e^{i\pi}+1=0 \label{eq:euler}
\end{equation}

\begin{align}
  (a+b)^2 &= a^2+2ab+b^2,\\
  (a-b)^2 &= a^2-2ab+b^2.
\end{align}
\end{lstlisting}

例如欧拉公式（\eqref{eq:euler}）被许多人称为数学中最美的公式之一：
\begin{equation}
  e^{i\pi}+1=0 \label{eq:euler}
\end{equation}

\subsection{插图}

插图既可以使用外部图片（\texttt{graphicx}），也可以直接用 \texttt{tikz}
绘制矢量图，后者无需额外文件、更利于保持文档自包含：

\begin{lstlisting}
\begin{figure}[htbp]
  \centering
  \begin{tikzpicture}
    \draw[fill=structurecolor!20] (0,0) circle (1);
    \draw[->,thick] (0,0) -- (0.707,0.707);
    \node at (0.4,0.9) {\scriptsize $\theta$};
  \end{tikzpicture}
  \caption{一个单位圆}\label{fig:unitcircle}
\end{figure}
\end{lstlisting}

引用插图时可以使用文档类提供的 \verb|\figref|（或
\verb|\tabref| 引用表格），例如"见图~\figref{fig:unitcircle}"：

\begin{figure}[htbp]
  \centering
  \begin{tikzpicture}
    \draw[fill=structurecolor!20] (0,0) circle (1);
    \draw[->,thick] (0,0) -- (0.707,0.707);
    \node at (0.42,0.55) {\scriptsize $\theta$};
    \draw[<->] (1.2,0) arc (0:45:1.2);
  \end{tikzpicture}
  \caption{一个单位圆}\label{fig:unitcircle}
\end{figure}

\subsection{表格}

三线表使用 \texttt{booktabs} 的 \verb|\toprule|、
\verb|\midrule|、\verb|\bottomrule|：

\begin{lstlisting}
\begin{table}[htbp]
  \centering
  \caption{示例数据}\label{tab:data}
  \begin{tabular}{lcc}
    \toprule
    变量 & 均值 & 方差 \\
    \midrule
    $x$ & 0.0 & 1.0 \\
    $y$ & 1.5 & 0.4 \\
    \bottomrule
  \end{tabular}
\end{table}
\end{lstlisting}

\subsection{代码列表}

源代码环境由内置的 \texttt{listings} 配置提供，
本手册中的全部示例代码均使用 \texttt{lstlisting} 环境排版。

\subsection{列表}

无序、有序、描述三种列表均可用，且项目符号被替换为主题色圆点：

\begin{lstlisting}
\begin{itemize}
  \item 无序列表项
\end{itemize}
\begin{enumerate}
  \item 有序列表项
\end{enumerate}
\begin{description}
  \item[术语] 解释文字。
\end{description}
\end{lstlisting}

\section{参考文献的使用}

参考文献数据库写在 \texttt{refs.bib} 中，主文件前言用
\verb|\addbibresource{refs.bib}| 引入：

\begin{lstlisting}
% refs.bib
@book{lamport1994latex,
  author    = {Lamport, Leslie},
  title     = {LaTeX: A Document Preparation System},
  edition   = {2nd},
  publisher = {Addison-Wesley},
  year      = {1994},
}
\end{lstlisting}

正文中用 \verb|\cite{...}| 引用，文末用
\verb|\printbibliography| 输出文献列表。例如
\cite{lamport1994latex} 是 LaTeX 的经典入门著作。注意：引入了
文献后，编译时必须在两次 \texttt{xelatex} 之间运行 \texttt{biber}，
否则引用位置会出现问号。

\section{常见问题与排错}

\begin{description}
  \item[\texttt{elegantbook.cls} 找不到] 说明发行版未安装该文档类。
        解决方法：运行 \texttt{tlmgr install elegantbook}（TeX~Live），
        或者把 \texttt{elegantbook.cls} 直接复制到源文件所在目录（推荐，
        本手册即采用此方式保证自包含）；
  \item[\texttt{File xxx.sty not found}] 缺少宏包，用
        \texttt{tlmgr install xxx} 安装，或检查是否误用了 pdfLaTeX
        编译中文文档；
  \item[编译后中文乱码/缺字] 确认源文件为 UTF-8 编码，并确认系统中
        存在中文字体（Windows 自带宋体/黑体/楷体）；
  \item[引用显示为问号] 交叉引用需多编一次；参考文献需运行
        \texttt{biber}；
  \item[\texttt{biber} 报错] 检查主文件选项 \texttt{bibend} 与后端
        是否一致（默认 \texttt{biber}，若用 BibTeX 后端则
        需 \texttt{bibend=bibtex}）。
\end{description}

\section{最小示例逐行讲解}

把第一节的最小示例逐行拆开，说明每一行的作用：

\begin{table}[htbp]
  \centering
  \caption{最小示例逐行注释}\label{tab:linebyline}
  \begin{tabular}{lp{9.5cm}}
    \toprule
    代码 & 作用 \\
    \midrule
    \texttt{\textbackslash documentclass[lang=cn]\{elegantbook\}}
      & 声明使用 ElegantBook 文档类，启用中文模式；
        之后才可编写正文 \\
    \texttt{\textbackslash title\{...\}} 等四行
      & 填写封面元信息：书名、作者、单位、日期 \\
    \texttt{\textbackslash begin\{document\}}
      & 正文开始（此前为导言区，只能放设置命令）\\
    \texttt{\textbackslash maketitle}
      & 按元信息自动排版整页封面并输出 \\
    \texttt{\textbackslash tableofcontents}
      & 输出目录（内容来自 \texttt{.toc} 文件）\\
    \texttt{\textbackslash chapter\{引言\}}
      & 开启一章；章标题自动进入目录与页眉 \\
    \texttt{\textbackslash end\{document\}}
      & 正文结束；\texttt{\textbackslash printbibliography} 等
        收尾命令应在它之前 \\
    \bottomrule
  \end{tabular}
\end{table}

注意导言区（\verb|\begin{document}| 之前）只能放\textbf{设置类}命令，
正文内容必须放在文档环境内——这是 LaTeX 最基本的语法边界。

\section{辅助文件详解}

编译过程中产生的辅助文件各有分工，了解它们有助于排查问题：

\begin{table}[htbp]
  \centering
  \caption{辅助文件一览}\label{tab:auxfiles}
  \begin{tabular}{llp{7.5cm}}
    \toprule
    文件 & 生产者 & 内容与作用 \\
    \midrule
    \texttt{.aux} & xelatex & 交叉引用表：\texttt{\textbackslash label}
          与 \texttt{\textbackslash ref} 的信息交换 \\
    \texttt{.toc} & xelatex & 目录条目，由 \texttt{\textbackslash
          chapter} 等写入 \\
    \texttt{.out} & hyperref & PDF 书签树 \\
    \texttt{.bcf} & biblatex & 文献引用记录（第一次 xelatex 写入）\\
    \texttt{.bbl} & biber & 文献列表排版结果，xelatex 直接读入 \\
    \texttt{.blg} & biber & biber 运行日志（报错时先看它）\\
    \texttt{.log} & xelatex & 编译日志（报错时的主要线索）\\
    \texttt{.synctex.gz} & xelatex & 源码与 PDF 双向定位（编辑器用）\\
    \bottomrule
  \end{tabular}
\end{table}

这些文件都可放心删除——下次编译会自动重建。唯一需要保留的是
\texttt{main.tex}、\texttt{chapters/} 与 \texttt{refs.bib}。

\section{三种编译方式}

\subsection{手动命令（最透明）}

适合理解原理与排查问题：
\begin{verbatim}
xelatex main.tex && biber main && xelatex main.tex && xelatex main.tex
\end{verbatim}

\subsection{latexmk（写作期推荐）}

latexmk 自动判断需要重跑几遍。在项目根目录放一个
\texttt{.latexmkrc}（Windows 上亦可），内容为：
\begin{lstlisting}
$pdf_mode     = 5;      # 用 xelatex 编译
$bibtex_use   = 2;      # 需要时自动运行参考文献工具
$biber        = 'biber %O %S';   # 指定 biber 后端
\end{lstlisting}
之后只需一条命令：
\begin{verbatim}
latexmk main.tex
\end{verbatim}

\subsection{编辑器集成}

\begin{description}
  \item[TeXstudio] 选项 $\rightarrow$ 构建 $\rightarrow$ 默认编译器
        选 \texttt{XeLaTeX}；文献工具选 \texttt{biber}；
  \item[VS Code + LaTeX Workshop] 在 \texttt{settings.json} 中把
        \texttt{latex-workshop.latex.recipes} 第一条配成
        \texttt{xelatex $\rightarrow$ biber $\rightarrow$ xelatex $\rightarrow$
        xelatex}；
  \item[Overleaf] 菜单 $\rightarrow$ 编译器选 \texttt{XeLaTeX}；
        Overleaf 会自动调用 biber。
\end{description}

\section{常用环境速查}

\subsection{数学公式环境}

文档类已加载 amsmath，常用环境如下：

\begin{table}[htbp]
  \centering
  \caption{常用数学环境}\label{tab:mathenv}
  \begin{tabular}{lll}
    \toprule
    环境 & 用途 & 编号 \\
    \midrule
    \texttt{equation} / \texttt{equation*} & 单行公式 & 有/无 \\
    \texttt{align} / \texttt{align*} & 多行按 \texttt{\&} 对齐 & 有/无 \\
    \texttt{gather} / \texttt{gather*} & 多行居中 & 有/无 \\
    \texttt{multline} & 超长公式换行（首行左、末行右）& 有 \\
    \texttt{cases} & 分段函数 & 嵌套使用 \\
    \texttt{split} & 多行共享一个编号 & 共享 \\
    \texttt{pmatrix}/\texttt{bmatrix} & 圆/方括号矩阵 & --- \\
    \bottomrule
  \end{tabular}
\end{table}

例如分段函数与矩阵：

\begin{lstlisting}
\[
  f(x)=\begin{cases}
    x^2, & x\geqslant 0,\\
    -x,  & x<0,
  \end{cases}
  \qquad
  A=\begin{pmatrix} a & b \\ c & d \end{pmatrix}
\]
\end{lstlisting}

渲染结果为
\[
  f(x)=\begin{cases}
    x^2, & x\geqslant 0,\\
    -x,  & x<0,
  \end{cases}
  \qquad
  A=\begin{pmatrix} a & b \\ c & d \end{pmatrix}.
\]

\subsection{其他常用环境}

\begin{description}
  \item[\texttt{itemize}/\texttt{enumerate}/\texttt{description}]
        无序/有序/描述列表，项目符号为主题色圆点；
  \item[\texttt{lstlisting}] 代码环境，支持语法高亮；
  \item[\texttt{verbatim}] 原样输出（不处理命令）；
  \item[\texttt{quote}/\texttt{quotation}] 短/长引文；
  \item[\texttt{abstract}] 摘要（文档类自带样式）；
  \item[\texttt{introduction}] 章前引言框（双栏要点列表）。
\end{description}

\section{标签与交叉引用规范}

\verb|\label| 与 \verb|\ref| 是文档内引用的基础。推荐的
标签命名规范是"类型前缀 + 冒号 + 语义名"：

\begin{table}[htbp]
  \centering
  \caption{标签前缀约定}\label{tab:labels}
  \begin{tabular}{lll}
    \toprule
    前缀 & 对象 & 引用命令 \\
    \midrule
    \texttt{fig:} & 图 & \verb|\ref|、\verb|\figref| \\
    \texttt{tab:} & 表 & \verb|\ref|、\verb|\tabref| \\
    \texttt{eq:} & 公式 & \verb|\eqref|（自动加括号）\\
    \texttt{sec:}/\texttt{ch:} & 节/章 & \verb|\ref|、\verb|\autoref| \\
    \texttt{thm:}/\texttt{def:}/\texttt{lem:} & 定理/定义/引理 & \verb|\ref| \\
    \texttt{alg:} & 算法 & \verb|\ref| \\
    \bottomrule
  \end{tabular}
\end{table}

\verb|\label| 必须放在产生编号的环境\emph{内部}（例如
\verb|\begin{equation}...\label{eq:x}\end{equation}|），
它记录的是"当前计数器的值"，放在外面会捕获错误编号。

\section{章节层级与编号深度}

ElegantBook 继承 \texttt{book} 类的层级结构：
\begin{equation*}
  \texttt{part} \succ \texttt{chapter} \succ \texttt{section}
  \succ \texttt{subsection} \succ \texttt{subsubsection}
  \succ \texttt{paragraph} \succ \texttt{subparagraph}.
\end{equation*}
默认编号深度 \texttt{secnumdepth=2}（即编号到
\texttt{subsection}），目录深度 \texttt{tocdepth=2}。
需要更深时可调整：
\begin{lstlisting}
\setcounter{secnumdepth}{3}   % subsubsection 也编号
\setcounter{tocdepth}{3}      % 目录包含 subsubsection
\end{lstlisting}

\section{一个完整的示例章}

综合以上内容，下面是一章完整的写法模板
（保存为 \texttt{chapters/ch01\_demo.tex} 并被主文件
\verb|\include|）：

\begin{lstlisting}
% chapters/ch01_demo.tex
\chapter{示例章}

\section{引入}
本节对应 \ref{sec:model} 的模型。

\begin{definition}[导数]\label{def:deriv}
设 $f$ 在 $x_0$ 附近有定义，若极限
\[ f'(x_0)=\lim_{h\to0}\frac{f(x_0+h)-f(x_0)}{h} \]
存在，则称 $f$ 在 $x_0$ 处可导。
\end{definition}

\begin{theorem}[可导必连续]
若 $f$ 在 $x_0$ 处可导，则 $f$ 在 $x_0$ 处连续。
\end{theorem}
\begin{proof}
由 $f'(x_0)$ 存在直接推出
$\lim_{x\to x_0}\big(f(x)-f(x_0)\big)=0$。
\end{proof}

\section{模型}\label{sec:model}
线性近似（\eqref{eq:linear}）是微积分的基本思想：
\begin{equation}
  f(x)\approx f(x_0)+f'(x_0)(x-x_0). \label{eq:linear}
\end{equation}

\begin{figure}[htbp]
  \centering
  \begin{tikzpicture}
    \draw[->] (-0.5,0) -- (3,0) node[right] {$x$};
    \draw[->] (0,-0.5) -- (0,3) node[above] {$y$};
    \draw[thick, structurecolor] plot[domain=0:2.5]
      (\x, {0.5*\x*\x});
    \draw[thick, second] (0.5,0.125) -- (2.5,2.125);
  \end{tikzpicture}
  \caption{割线逼近切线}\label{fig:approx}
\end{figure}

图像见图~\figref{fig:approx}，理论见 \cite{knuth1984texbook}。
\end{lstlisting}

按此模板组织，一本书的每个章节彼此独立，便于多人协作与单独编译调试。

...
\end{lstlisting}

\verb|\includeonly| 与 \verb|\include| 的配对使用是 LaTeX
标准做法：未包含的章节不重新排版，但\emph{交叉引用与页码
保持不变}（.aux 中保留了它们的记录）。配合 git 等版本控制，
多作者可以按章并行写作、由一人负责主文件集成。
